\textbf{SIMULAÇÃO ESPACIAL DA TRANSMISSÃO DA DENGUE NO MUNICÍPIO DE SÃO
PAULO POR MODELAGEM BASEADA EM AGENTES NA PLATAFORMA GAMA}

Texto acadêmico consolidado --- versão de revisão técnica \textbar{} 8
nov. 2026

Autoria, curso, instituição, orientador e demais elementos pré-textuais:
preencher conforme o modelo institucional efetivamente adotado no
repositório. Este documento não substitui a capa, a ficha catalográfica
nem a folha de aprovação do template institucional.

\section{RESUMO}\label{resumo}

Este trabalho apresenta a documentação e a avaliação em desenvolvimento
de um simulador espacial e estocástico da transmissão da dengue no
município de São Paulo, implementado na plataforma GAMA por meio da
linguagem GAML. A abordagem baseada em agentes representa indivíduos
humanos, mosquitos vetores e componentes ambientais, contemplando
estados epidemiológicos, deslocamentos, interações por proximidade,
informações geográficas, condições climáticas e registros temporais. A
metodologia compreende análise do código-fonte, especificação dos
comportamentos, verificação de invariantes, planejamento de experimentos
controlados e avaliação de reprodutibilidade. No repositório foram
identificadas configurações de experimentos em lote para variabilidade
estocástica, vacinação, população vetorial e escalabilidade. Um conjunto
histórico de registros CSV permitiu identificar a necessidade de
conferir a correspondência entre totais populacionais e estados
epidemiológicos. Em 7 de novembro de 2026, o histórico do projeto
registrou modificações relacionadas a sementes, população de mosquitos e
inicialização vetorial; a eficácia dessas correções ainda demanda testes
de execução e auditoria dos resultados associados. Assim, os resultados
disponíveis sustentam uma análise de implementação e consistência, mas
não constituem validação epidemiológica ou previsão de surtos reais. O
estudo contribui para a explicitação de requisitos técnicos,
experimentais e de documentação necessários à avaliação de simuladores
baseados em agentes.

Palavras-chave: dengue; modelagem baseada em agentes; GAMA; GAML;
sistemas de informação geográfica; verificação de simulação.

\section{ABSTRACT}\label{abstract}

This study presents the documentation and ongoing assessment of a
spatial stochastic agent-based dengue transmission simulator for the
municipality of São Paulo, implemented in the GAMA platform using GAML.
The model represents human individuals, mosquito vectors, and
environmental components, including epidemiological states, mobility,
proximity-based interactions, geographic information, climatic
conditions, and time-series recording. The methodological approach
includes source-code analysis, behavioral specification, invariant
checking, controlled experimental design, and reproducibility
assessment. The repository contains batch experiment configurations
addressing stochastic variability, vaccination, vector population, and
scalability. An earlier CSV dataset revealed the need to audit
consistency between total populations and epidemiological state counts.
On November 7, 2026, the project history reported changes concerning
random seeds, mosquito population initialization, and vector states; the
effectiveness of these modifications still requires execution tests and
verification of associated outputs. Available evidence therefore
supports an implementation and consistency analysis but does not
establish epidemiological validation or outbreak forecasting capability.
The work contributes technical, experimental, and documentation
requirements for the rigorous assessment of agent-based simulators.

Keywords: dengue; agent-based modeling; GAMA; GAML; geographic
information systems; simulation verification.

Nota editorial: atualizar Resumo e Abstract somente após gerar e validar
as séries E03/E04/E05. Não inserir resultados quantitativos
inexistentes.

\section{SUMÁRIO --- ESTRUTURA}\label{sumuxe1rio-estrutura}

1 INTRODUÇÃO

2 FUNDAMENTAÇÃO TEÓRICA E TRABALHOS RELACIONADOS

3 MODELAGEM COMPUTACIONAL E JUSTIFICATIVA DA PLATAFORMA

4 ARQUITETURA E IMPLEMENTAÇÃO DO SIMULADOR

5 VERIFICAÇÃO, EXPERIMENTAÇÃO, RESULTADOS E DISCUSSÃO

6 CONCLUSÃO E TRABALHOS FUTUROS

REFERÊNCIAS

APÊNDICE A --- RASTREABILIDADE DAS EXIGÊNCIAS DO ORIENTADOR

A paginação definitiva e o sumário automático devem ser gerados no
template LaTeX/Senac; esta lista apresenta somente a estrutura do texto
consolidado.

\section{1 INTRODUÇÃO}\label{introduuxe7uxe3o}

A dengue é uma arbovirose causada pelo vírus dengue (DENV), cuja
transmissão ocorre predominantemente pela picada de fêmeas infectadas de
mosquitos do gênero Aedes. No Brasil, o Aedes aegypti constitui o
principal vetor urbano. São reconhecidos quatro sorotipos virais,
denominados DENV-1, DENV-2, DENV-3 e DENV-4, cuja circulação representa
um desafio à vigilância e à prevenção da doença (BRASIL, s.d.; WORLD
HEALTH ORGANIZATION, 2025). (Brasil, {[}s. d.{]}; World Health
Organization, 2025)

A infecção pode permanecer assintomática, provocar manifestações febris
ou evoluir para formas graves. A dinâmica de transmissão depende de
relações entre hospedeiros humanos, mosquitos vetores e características
do ambiente. O crescimento urbano, a disponibilidade de recipientes
favoráveis ao desenvolvimento do mosquito, as condições climáticas e a
movimentação das pessoas constituem fatores relevantes para compreender
a heterogeneidade espacial e temporal da transmissão (WORLD HEALTH
ORGANIZATION, 2025). (Brasil, {[}s. d.{]}; World Health Organization,
2025)

Segundo a Organização Mundial da Saúde (2025), aproximadamente metade da
população mundial encontra-se em áreas com risco de dengue. Embora
existam medidas preventivas e estratégias de imunização para públicos e
condições específicos, a prevenção depende fortemente do controle do
vetor, da redução dos criadouros, da vigilância e do acesso oportuno à
assistência em saúde. A relevância do problema motiva o uso de
abordagens computacionais capazes de explorar mecanismos e cenários sob
hipóteses explícitas. (Brasil, {[}s. d.{]}; World Health Organization,
2025)

O município de São Paulo oferece um contexto territorial complexo para
esse tipo de investigação. A concentração populacional, a diversidade de
usos do solo e a presença de deslocamentos entre diferentes regiões
justificam a investigação de uma representação espacial, desde que as
bases geográficas e demográficas sejam adequadamente identificadas.
Neste trabalho, a área de estudo corresponde ao município de São Paulo,
não apenas ao Distrito ou à Subprefeitura de Santo Amaro. (Brasil, {[}s.
d.{]}; World Health Organization, 2025)

A representação da dengue por simulação computacional exige escolhas que
transcendem a dimensão epidemiológica. É necessário definir entidades,
variáveis de estado, regras de execução, agendamento temporal, relações
de vizinhança e mecanismos de geração pseudoaleatória, além de preparar
dados geográficos e preservar a rastreabilidade das saídas. Essas
decisões afetam simultaneamente a interpretação do fenômeno, a
reprodutibilidade e o custo computacional do programa. (Macal; North,
2010)

A modelagem baseada em agentes --- Agent-Based Modeling (ABM) ---
permite representar indivíduos heterogêneos com atributos, estados e
comportamentos próprios. Suas interações com outros agentes e com o
ambiente podem gerar padrões coletivos que não são especificados
diretamente como saídas agregadas. Ferramentas como a plataforma GAMA
oferecem estruturas voltadas à implementação de modelos espacialmente
explícitos e à execução de experimentos computacionais (MACAL; NORTH,
2010; TAILLANDIER et al., 2019).

No campo de aplicação da dengue, Garcias (2024) desenvolveu um modelo na
plataforma GAMA para Portugal continental, com ênfase nas interações
entre humanos, mosquitos Aedes albopictus e ambiente. Oliveira et al.
(2025) apresentaram uma simulação multiagentes aplicada a Alto Santo,
Ceará, investigando parâmetros vetoriais e intervenções de controle.
Ambos os estudos demonstram possibilidades metodológicas, mas suas
espécies vetoriais, delimitações geográficas, parâmetros e resultados
não podem ser automaticamente transportados para uma aplicação em São
Paulo.

Neste contexto, o presente trabalho tem como objeto central a engenharia
e a análise de um simulador espacial de transmissão da dengue
desenvolvido na plataforma GAMA, por meio da linguagem GAML. A
implementação é examinada a partir do repositório público TccAbm,
observando a organização dos componentes, a inicialização dos agentes,
as regras de transmissão, a mobilidade, os parâmetros de vacinação, o
controle temporal, o ambiente geográfico e a produção de indicadores. A
leitura do código, por si só, não constitui evidência de funcionamento
correto ou validação epidemiológica, de modo que a verificação e a
avaliação experimental são tratadas como atividades específicas.
(Taillandier et al., 2019; Macal; North, 2010)

\subsection{1.1 Problema de pesquisa}\label{problema-de-pesquisa}

Um simulador pode gerar gráficos e séries temporais mesmo quando contém
transições de estados inconsistentes, problemas de escala temporal,
pressupostos espaciais insuficientemente documentados ou mecanismos
pseudoaleatórios de difícil reprodução. Dessa forma, a execução do
programa e a plausibilidade visual de suas saídas não são suficientes
para demonstrar a conformidade entre modelo conceitual, implementação e
propósito experimental. (HKIOSHI, 2026)

A análise preliminar da versão estudada do repositório TccAbm
identificou componentes separados para pessoas, mosquitos, ambiente,
criadouros e procedimentos globais. Também evidenciou a necessidade de
conferir, entre outros pontos, os estados epidemiológicos inicializados,
a conversão de escalas temporais, o controle das sementes e a
correspondência entre as geometrias carregadas e a área municipal. Esses
achados serão examinados detalhadamente nos capítulos de implementação e
avaliação; não são tratados, nesta introdução, como falhas confirmadas
por testes de execução. (HKIOSHI, 2026)

Diante desses desafios, estabelece-se a seguinte pergunta de pesquisa:
como desenvolver, documentar e avaliar computacionalmente um simulador
espacial e estocástico baseado em agentes na plataforma GAMA, capaz de
representar as interações entre humanos, mosquitos Aedes aegypti e
ambiente na dinâmica de transmissão da dengue no município de São Paulo,
considerando a consistência das regras, a reprodutibilidade e o
desempenho da implementação? (Taillandier et al., 2019; Macal; North,
2010)

\subsection{1.2 Hipótese de trabalho e
pressupostos}\label{hipuxf3tese-de-trabalho-e-pressupostos}

Parte-se do pressuposto de que a explicitação das entidades, transições,
escalas, fontes de dados e mecanismos de execução permite estabelecer
critérios verificáveis para analisar a implementação. Espera-se que a
relação sistemática entre especificação, código-fonte e testes revele
limites e propriedades do simulador que não podem ser inferidos apenas
de sua visualização. Trata-se de uma hipótese metodológica que deverá
ser examinada pelos procedimentos apresentados no Capítulo 5. (HKIOSHI,
2026)

A simulação não será considerada um mecanismo de previsão epidemiológica
validado sem comparação adequada com dados observados. A adoção de
população sintética reduzida, parâmetros hipotéticos ou agregação
territorial requer descrição explícita e impede extrapolações
automáticas para toda a população do município. (Macal; North, 2010)

\subsection{1.3 Justificativa}\label{justificativa}

A contribuição científica pretendida situa-se principalmente na Ciência
da Computação. Do ponto de vista de engenharia de software, a pesquisa
contempla modularização de código, organização de responsabilidades,
especificação das regras, rastreabilidade entre requisitos e
implementação, e verificação funcional. Do ponto de vista algorítmico,
envolve eventos probabilísticos, atualização de estados, consultas
espaciais e agendamento de processos. Na avaliação experimental,
interessa caracterizar variabilidade, sensibilidade e desempenho sob
configurações controladas. (HKIOSHI, 2026)

A escolha do GAMA se justifica por suas capacidades de representar
agentes espacialmente explícitos, incorporar dados territoriais e
organizar experimentos. A plataforma constitui o ambiente computacional
da implementação, mas sua adoção não garante, isoladamente, que as
regras biológicas, as operações espaciais ou o desempenho estejam
corretos (TAILLANDIER et al., 2019).

A utilidade da dengue como domínio de aplicação está em oferecer um
sistema com interações locais, agentes heterogêneos, eventos
estocásticos e variáveis ambientais. Assim, o trabalho combina um
problema socialmente relevante com desafios técnicos de modelagem e
implementação. Sua contribuição deverá ser demonstrada por documentação
e evidências, e não apenas pela escolha de uma nova cidade ou pela
reprodução de funcionalidades descritas na literatura. (HKIOSHI, 2026)

\subsection{1.4 Objetivos}\label{objetivos}

1.4.1 Objetivo geral

Desenvolver, documentar e avaliar um simulador computacional espacial,
estocástico e baseado em agentes, utilizando a plataforma GAMA e dados
de Sistemas de Informação Geográfica (SIG), para representar a dinâmica
de transmissão da dengue no município de São Paulo, por meio das
interações entre agentes humanos, mosquitos Aedes aegypti e componentes
ambientais, com ênfase na modelagem e na implementação dos algoritmos,
na verificação funcional, na reprodutibilidade e na análise do
desempenho computacional.

1.4.2 Objetivos específicos

1. Investigar os fundamentos da transmissão da dengue e da modelagem
baseada em agentes, identificando requisitos biológicos, espaciais e
computacionais necessários à especificação do simulador.

2. Especificar e documentar os agentes, atributos, estados
epidemiológicos, comportamentos, interações e hipóteses de representação
adotados no modelo.

3. Justificar a utilização de GAMA/GAML e sua integração com SIG,
comparando-a qualitativamente a alternativas de simulação conforme os
requisitos do projeto.

4. Analisar e documentar a implementação dos módulos responsáveis por
humanos, mosquitos Aedes aegypti, ambientes, criadouros e procedimentos
globais, distinguindo funcionalidades presentes de aperfeiçoamentos
propostos.

5. Caracterizar e verificar a inicialização, a distribuição espacial, a
mobilidade, a vacinação, o controle temporal, os eventos probabilísticos
e os mecanismos de geração de saídas.

6. Definir e aplicar, conforme a disponibilidade de execução, testes de
consistência de estados, transições, parâmetros e indicadores,
registrando as evidências e as limitações encontradas.

7. Planejar e executar, quando tecnicamente viável, experimentos com
diferentes sementes e parâmetros, de modo a examinar a variabilidade
estocástica e a sensibilidade dos resultados.

8. Avaliar, mediante medições reproduzíveis, o custo computacional em
diferentes configurações de agentes, incluindo tempo de processamento e
consumo de memória.

9. Investigar a compatibilidade entre indicadores simulados e séries
epidemiológicas oficiais do município de São Paulo, realizando
comparação quantitativa apenas quando as escalas, períodos e definições
permitirem.

10. Sistematizar as limitações identificadas na implementação e nos
dados, propondo aperfeiçoamentos técnicos voltados à confiabilidade, à
reprodutibilidade e à capacidade de experimentação.

\subsection{1.5 Delimitação e escopo}\label{delimitauxe7uxe3o-e-escopo}

O recorte territorial do estudo é o município de São Paulo, Estado de
São Paulo, Brasil. Sua representação computacional será descrita a
partir das bases efetivamente utilizadas pela implementação. O
carregamento de geometrias de subprefeituras ou a utilização de regiões
agregadas não equivale, por si só, à reprodução completa da população,
da mobilidade e da distribuição real do risco epidemiológico em todos os
distritos da cidade. (HKIOSHI, 2026)

O domínio epidemiológico abrange a representação da transmissão
envolvendo humanos e mosquitos Aedes aegypti, conforme as regras
adotadas no simulador sob estudo. Os estados compartimentais e os
componentes ambientais serão apresentados em duas camadas: especificação
conceitual desejada e comportamento efetivamente encontrado no código.
Diferenças entre ambas serão registradas como objetos de análise e,
quando possível, de verificação. (HKIOSHI, 2026)

Este trabalho não pressupõe representação simultânea dos quatro
sorotipos, reprodução integral das características imunológicas da
dengue, equivalência demográfica perfeita ou precisão preditiva de casos
futuros. A eventual calibração ou validação com dados oficiais dependerá
da disponibilidade de séries com período, área e indicadores
compatíveis, e de uma regra que diferencie infecções simuladas de
notificações observadas. (Grimm et al., 2020; Taillandier et al., 2019)

\subsection{1.6 Procedimentos
metodológicos}\label{procedimentos-metodoluxf3gicos}

O trabalho é conduzido como estudo de desenvolvimento, documentação e
avaliação computacional de uma implementação disponível no repositório
TccAbm. Inicialmente, realizou-se revisão bibliográfica sobre dengue,
modelagem baseada em agentes e plataformas espaciais. Em seguida, foram
examinados a estrutura do código GAML e os dados geográficos e
experimentais disponíveis, com o propósito de distinguir os mecanismos
implementados das funcionalidades apenas propostas. A autoria, a
participação dos desenvolvedores e o histórico das alterações deverão
ser apresentados de acordo com as evidências documentais disponíveis.
(HKIOSHI, 2026)

A especificação do modelo utiliza o protocolo ODD como referência de
documentação (GRIMM et al., 2020). A avaliação computacional é
organizada em verificação funcional, reprodutibilidade, sensibilidade e
desempenho. Testes e medições somente serão relatados como realizados
quando acompanhados de configuração, procedimento, artefatos de saída e
resultados verificáveis. A comparação com dados observados de São Paulo,
quando possível, será tratada separadamente da verificação do software.

Os registros já existentes permitem uma análise exploratória e a
identificação de verificações necessárias, mas não demonstram,
isoladamente, validação epidemiológica nem escalabilidade. O Capítulo 5
separa evidências disponíveis, resultados de experimentos executados e
protocolos experimentais pendentes, sem confundir planejamento com
resultados observados. (Grimm et al., 2020; Taillandier et al., 2019)

\subsection{1.7 Organização do
trabalho}\label{organizauxe7uxe3o-do-trabalho}

Além desta introdução, o Capítulo 2 apresenta a fundamentação
epidemiológica e computacional e compara trabalhos relacionados. O
Capítulo 3 especifica o modelo, as entidades, as escalas, os mecanismos
e as hipóteses. O Capítulo 4 examina a arquitetura do código GAML, a
inicialização, o processamento dos comportamentos, a integração GIS e a
coleta dos indicadores. O Capítulo 5 reúne os procedimentos de
verificação, as evidências experimentais efetivamente obtidas, as
limitações e a avaliação de reprodutibilidade e desempenho. O Capítulo 6
sintetiza as contribuições demonstradas e propõe melhorias e trabalhos
futuros. (HKIOSHI, 2026)

\section{2 FUNDAMENTAÇÃO TEÓRICA E TRABALHOS
RELACIONADOS}\label{fundamentauxe7uxe3o-teuxf3rica-e-trabalhos-relacionados}

Este capítulo reúne os fundamentos necessários à compreensão,
especificação e avaliação de um simulador espacial baseado em agentes
para a transmissão da dengue no município de São Paulo. Inicialmente,
são apresentados os mecanismos epidemiológicos indispensáveis à
definição das entidades e dos estados do sistema. Em seguida, são
discutidas as abordagens de modelagem, os conceitos computacionais, a
integração de dados geográficos, os mecanismos de estocasticidade e a
plataforma GAMA. Por fim, os trabalhos relacionados são examinados
criticamente, distinguindo suas contribuições da implementação
investigada neste TCC. (Brasil, {[}s. d.{]}; World Health Organization,
2025)

A revisão mantém uma separação metodológica: a literatura epidemiológica
fundamenta as hipóteses do fenômeno; a literatura de simulação e
engenharia de software fundamenta a representação, a implementação e a
avaliação do artefato computacional. Nenhum resultado relatado nos
estudos anteriores é atribuído ao simulador desta pesquisa.

\subsection{2.1 Fundamentos epidemiológicos da
dengue}\label{fundamentos-epidemioluxf3gicos-da-dengue}

\subsubsection{2.1.1 Doença, agente etiológico e
sorotipos}\label{doenuxe7a-agente-etioluxf3gico-e-sorotipos}

A dengue é uma arbovirose causada pelo vírus da dengue, transmitido
principalmente por mosquitos do gênero Aedes. Quatro sorotipos,
designados DENV-1, DENV-2, DENV-3 e DENV-4, participam da dinâmica
epidemiológica. A infecção pode ser assintomática ou apresentar
manifestações clínicas de gravidade variável. Infecções posteriores por
sorotipos diferentes exigem atenção especial à interação entre imunidade
prévia e risco de doença grave (WORLD HEALTH ORGANIZATION, 2026).
(Brasil, {[}s. d.{]}; World Health Organization, 2025)

Para fins de modelagem computacional, a existência de diferentes
sorotipos demonstra que um compartimento genérico de recuperação não
equivale automaticamente à imunidade definitiva contra todas as
infecções futuras. Uma representação com apenas um sorotipo constitui
simplificação legítima quando explicitada, mas limita a interpretação
epidemiológica das saídas. Essa restrição deverá ser considerada na
definição dos estados e na leitura dos resultados do simulador. (Brasil,
{[}s. d.{]}; World Health Organization, 2025)

\subsubsection{2.1.2 Vetor e ciclo de
desenvolvimento}\label{vetor-e-ciclo-de-desenvolvimento}

No contexto urbano brasileiro, Aedes aegypti é o vetor de maior
relevância para o modelo estudado. A transmissão depende das fêmeas
adultas, enquanto os estágios de ovo, larva e pupa compõem o
desenvolvimento pré-adulto. A temperatura, a disponibilidade de água e
as características do ambiente influenciam diferentes etapas do ciclo;
entretanto, a presença de água ou vegetação não constitui, isoladamente,
uma medição direta da abundância vetorial. (Brasil, {[}s. d.{]}; World
Health Organization, 2025)

A distinção entre desenvolvimento biológico e transmissão é importante
do ponto de vista da engenharia: o sistema pode representar indivíduos
imaturos de modo agregado, sem que isso implique que todos os estágios
sejam objetos individuais. O nível de detalhamento deve ser justificado
em relação à pergunta de pesquisa, à disponibilidade de dados e ao custo
computacional. Garcias (2024) discute comportamentos de Aedes albopictus
em Portugal, mas seus parâmetros não podem ser transferidos sem exame
para Aedes aegypti em São Paulo.

\subsubsection{2.1.3 Transmissão humano--vetor e períodos de
incubação}\label{transmissuxe3o-humanovetor-e-peruxedodos-de-incubauxe7uxe3o}

A transmissão envolve interações entre humanos e mosquitos. Uma fêmea
suscetível pode adquirir o vírus ao realizar alimentação sanguínea em um
humano infectante. Após o período de incubação extrínseca, o vetor pode
tornar-se infectante. Um humano suscetível pode adquirir o vírus quando
recebe a picada de uma fêmea infectante. Essas etapas dependem de
condições biológicas e da oportunidade de contato, não apenas da
proximidade geométrica. (Brasil, {[}s. d.{]}; World Health Organization,
2025)

Na modelagem, é necessário distinguir infecção, infectividade e
manifestação clínica. Em particular, nem todo indivíduo recém-infectado
é imediatamente capaz de transmitir o vírus; por isso a fase exposta
requer definição cuidadosa. Uma inconsistência entre a descrição
conceitual e a condição implementada no código deve ser tratada como
questão de verificação, e não como prova de que o modelo representa o
mecanismo biológico desejado. (Brasil, {[}s. d.{]}; World Health
Organization, 2025)

\subsubsection{2.1.4 Fatores ambientais, urbanos e medidas de
controle}\label{fatores-ambientais-urbanos-e-medidas-de-controle}

A dinâmica da dengue depende simultaneamente de fatores biológicos,
climáticos e sociais. Temperatura, precipitação, disponibilidade de
criadouros, distribuição das pessoas e deslocamentos urbanos podem
modificar as oportunidades de encontro entre hospedeiros e vetores. No
município de São Paulo, a heterogeneidade territorial torna importante
distinguir limites administrativos, unidades de agregação, densidade
populacional e escala de representação. (Brasil, {[}s. d.{]}; World
Health Organization, 2025)

As medidas de prevenção e controle podem incluir eliminação ou manejo de
criadouros, redução da exposição a picadas, intervenções vetoriais e
vacinação segundo critérios de saúde pública. Em modelos computacionais,
cada intervenção exige um mecanismo operacional definido. Por exemplo,
reduzir a probabilidade de infecção não é equivalente a reduzir a
população de mosquitos; eficiência contra doença sintomática tampouco
deve ser automaticamente utilizada como proteção contra infecção.
(Brasil, {[}s. d.{]}; World Health Organization, 2025)

\subsection{2.2 Modelagem e simulação
computacional}\label{modelagem-e-simulauxe7uxe3o-computacional}

Modelar significa produzir uma representação deliberadamente
simplificada de um sistema real para responder a perguntas determinadas.
A simulação é a execução temporal dessas regras por um programa. Deve-se
diferenciar o fenômeno real, o modelo conceitual, a implementação em
software e os dados de saída: uma falha pode decorrer de hipóteses
inadequadas, erros de implementação, dados incompatíveis ou problemas
experimentais. (Macal; North, 2010; Hinsch, {[}s. d.{]})

\subsubsection{2.2.1 Abordagens de
simulação}\label{abordagens-de-simulauxe7uxe3o}

A dinâmica de sistemas representa relações agregadas por estoques e
fluxos; a simulação a eventos discretos descreve alterações de estado
disparadas por eventos; e a modelagem baseada em agentes representa
entidades individuais, seus atributos e suas interações. As abordagens
não são mutuamente exclusivas: um modelo de agentes pode usar estados
compartimentais e produzir indicadores agregados ao longo do tempo
(MACAL; NORTH, 2010). (Macal; North, 2010; Hinsch, {[}s. d.{]})

Quadro 1 --- Abordagens computacionais de simulação

\begin{longtable}[]{@{}
  >{\raggedright\arraybackslash}p{(\linewidth - 4\tabcolsep) * \real{0.3333}}
  >{\raggedright\arraybackslash}p{(\linewidth - 4\tabcolsep) * \real{0.3333}}
  >{\raggedright\arraybackslash}p{(\linewidth - 4\tabcolsep) * \real{0.3333}}@{}}
\toprule\noalign{}
\begin{minipage}[b]{\linewidth}\raggedright
\textbf{Abordagem}
\end{minipage} & \begin{minipage}[b]{\linewidth}\raggedright
\textbf{Unidade predominante}
\end{minipage} & \begin{minipage}[b]{\linewidth}\raggedright
\textbf{Limitação relevante ao problema}
\end{minipage} \\
\midrule\noalign{}
\endhead
\bottomrule\noalign{}
\endlastfoot
Dinâmica de sistemas & Variáveis agregadas & Pode ocultar
heterogeneidade espacial individual \\
Eventos discretos & Eventos e recursos & Exige construção explícita das
entidades e relações geográficas \\
Baseada em agentes & Entidades e interações locais & Aumenta demandas de
especificação, teste e processamento \\
\end{longtable}

Fonte: elaboração própria, a partir das referências citadas.

\subsubsection{2.2.2 Modelos preditivos e
explicativos}\label{modelos-preditivos-e-explicativos}

A finalidade de um modelo condiciona as abstrações utilizadas, os dados
necessários e os critérios pelos quais seus resultados devem ser
avaliados. Hinsch (s.d.), em apresentação introdutória sobre modelagem,
distingue duas orientações: a preditiva, voltada à estimação do
comportamento de um sistema sob condições especificadas, e a
explicativa, voltada à investigação dos mecanismos capazes de gerar
determinados padrões observáveis.

Na modelagem preditiva, importa examinar a qualidade das estimativas
diante de observações independentes e de situações que não participaram
do ajuste do modelo. Na modelagem explicativa, o interesse concentra-se
na relação entre hipóteses, regras e padrões emergentes: se uma
combinação de mecanismos reproduz certo comportamento, isso fornece
evidência de sua suficiência dentro do modelo, mas não demonstra que
seja a única explicação possível para o fenômeno (HINSCH, s.d.). (Macal;
North, 2010; Hinsch, {[}s. d.{]})

Essa distinção não constitui uma separação absoluta entre tipos de
software: um mesmo simulador pode sustentar investigações explicativas
e, após calibração e avaliação apropriadas, ser utilizado em tarefas
preditivas. Entretanto, resultados visualmente semelhantes a séries
observadas não comprovam, por si só, capacidade de previsão. A
interpretação das saídas depende da finalidade declarada, da qualidade
das bases de entrada, das hipóteses de representação e do protocolo
experimental. (Macal; North, 2010; Hinsch, {[}s. d.{]})

No presente TCC, o simulador de dengue aplicado ao município de São
Paulo é tratado predominantemente como instrumento explicativo e
exploratório. Sua utilização permite investigar as consequências das
regras computacionais de mobilidade, contato humano--vetor, evolução
epidemiológica e influência ambiental. Essa classificação não significa
que o modelo real já foi validado: a verificação funcional, a análise de
sensibilidade, a reprodutibilidade e eventual confronto com dados
epidemiológicos observados constituem etapas distintas que deverão ser
documentadas no Capítulo 5. (Macal; North, 2010; Hinsch, {[}s. d.{]})

O caráter explicativo orienta também a seleção de detalhes do modelo.
Incluir novos estados, agentes ou mecanismos somente contribui para a
pesquisa quando esses elementos são necessários às perguntas
investigadas ou permitem testar hipóteses relevantes. Portanto, as
simplificações devem ser justificadas e suas consequências avaliadas, em
vez de presumir que maior complexidade necessariamente produz maior
adequação (HINSCH, s.d.). (Macal; North, 2010; Hinsch, {[}s. d.{]})

\subsection{2.3 Modelagem e simulação baseada em
agentes}\label{modelagem-e-simulauxe7uxe3o-baseada-em-agentes}

A modelagem baseada em agentes (Agent-Based Modeling --- ABM) descreve
sistemas por agentes com atributos, estados e comportamentos. A
simulação executa as regras dos agentes em determinado ambiente,
permitindo observar padrões coletivos produzidos por interações locais.
Macal e North (2010) destacam a representação de entidades e relações
como característica central dessa abordagem. (Santos, A. R.; Santos, F.,
2024)

Para a dengue, a abordagem permite manter separados atributos
individuais, como localização e condição vacinal, e estados dinâmicos,
como suscetibilidade e infecção. É possível estabelecer diferentes
oportunidades de contato dependendo do posicionamento e da mobilidade,
sem assumir que todos os indivíduos se misturam homogeneamente. Essa
flexibilidade, porém, produz desafios de programação, inicialização e
calibração. (Macal; North, 2010; Taillandier et al., 2019)

\subsubsection{2.3.1 Entidades, atributos, comportamentos e
emergência}\label{entidades-atributos-comportamentos-e-emerguxeancia}

Um agente pode ser descrito por identificador, tipo, atributos
constantes, variáveis de estado e ações. Comportamentos são operações
que consultam o estado do agente e o ambiente e podem modificá-los. A
emergência designa padrões agregados gerados por essas interações, como
concentrações espaciais ou curvas temporais; não deve ser confundida com
validade empírica automática do modelo. (Macal; North, 2010; Taillandier
et al., 2019)

\subsubsection{2.3.2 Vizinhanças e custo das
interações}\label{vizinhanuxe7as-e-custo-das-interauxe7uxf5es}

Em simulações espaciais, uma interação depende de uma regra de
vizinhança que identifica possíveis parceiros. A busca exaustiva entre
N\_H humanos e N\_M mosquitos envolve até N\_H × N\_M comparações por
etapa. Índices e consultas espaciais podem reduzir candidatos, mas seu
desempenho efetivo depende da distribuição dos agentes, da atualização
das posições e da implementação. A eficiência deve ser medida, não
presumida. (Macal; North, 2010; Taillandier et al., 2019)

\subsection{2.4 Modelos epidemiológicos compartimentais e máquinas de
estados}\label{modelos-epidemioluxf3gicos-compartimentais-e-muxe1quinas-de-estados}

Os modelos SIR, SEIR e SEI organizam fases epidemiológicas por estados.
Em uma representação individual, esses compartimentos podem ser
implementados como máquinas de estados: humanos suscetíveis (S),
expostos (E), infectantes (I) e recuperados (R), e mosquitos adultos
suscetíveis (S), expostos (E) e infectantes (I). A passagem entre
estados depende de eventos, temporizadores e probabilidades, conforme a
especificação adotada. (Oliveira et al., 2025; Garcias, 2024)

Essas siglas não garantem que o código execute corretamente as
transições. Uma máquina de estados deve definir o conjunto de estados
válidos, as condições de entrada e saída, as transições proibidas e o
comportamento terminal, como a morte. Para mosquitos, a transição
automática de infectante a suscetível não corresponde à estrutura SEI
habitualmente pretendida, razão pela qual necessita de justificativa ou
correção no componente implementado. (Oliveira et al., 2025; Garcias,
2024)

Quadro 2 --- Estados epidemiológicos e verificações necessárias

\begin{longtable}[]{@{}
  >{\raggedright\arraybackslash}p{(\linewidth - 4\tabcolsep) * \real{0.3333}}
  >{\raggedright\arraybackslash}p{(\linewidth - 4\tabcolsep) * \real{0.3333}}
  >{\raggedright\arraybackslash}p{(\linewidth - 4\tabcolsep) * \real{0.3333}}@{}}
\toprule\noalign{}
\begin{minipage}[b]{\linewidth}\raggedright
\textbf{Entidade}
\end{minipage} & \begin{minipage}[b]{\linewidth}\raggedright
\textbf{Estados conceituais}
\end{minipage} & \begin{minipage}[b]{\linewidth}\raggedright
\textbf{Ponto de verificação}
\end{minipage} \\
\midrule\noalign{}
\endhead
\bottomrule\noalign{}
\endlastfoot
Humano & S → E → I → R & Incubação, recuperação e infectividade \\
Mosquito adulto & S → E → I & Incubação extrínseca e permanência no
estado I \\
Vacinação & Atributo independente & Efeito sobre cada desfecho
epidemiológico \\
\end{longtable}

Fonte: elaboração própria, a partir das referências citadas.

\subsection{2.5 Representação espacial e Sistemas de Informação
Geográfica}\label{representauxe7uxe3o-espacial-e-sistemas-de-informauxe7uxe3o-geogruxe1fica}

Sistemas de Informação Geográfica (SIG) permitem organizar feições
territoriais e atributos localizados. Dados vetoriais podem representar
polígonos administrativos, vias e pontos de interesse; dados matriciais
podem representar superfícies ambientais. Em uma simulação espacial,
geometrias cumprem funções operacionais --- delimitar o espaço,
distribuir agentes, apoiar deslocamentos e localizar contatos --- e não
apenas ilustrativas. (Taillandier; Drogoul, 2010; Taillandier et al.,
2019)

A integração entre bases exige sistema de referência de coordenadas
coerente, unidades adequadas para cálculos de distância, geometrias
válidas e tratamento explícito de fronteiras e agregações. No presente
projeto, a área de estudo corresponde ao município de São Paulo. A
divisão em Centro, Norte, Leste, Sul e Oeste encontrada na implementação
deve ser descrita como uma forma de agregação do código, não confundida
automaticamente com divisões oficiais equivalentes ou com resolução
homogênea do território. (Taillandier; Drogoul, 2010; Taillandier et
al., 2019)

\subsubsection{2.5.1 Mobilidade humana e
redes}\label{mobilidade-humana-e-redes}

Mobilidade pode envolver seleção de destino e cálculo da trajetória,
procedimentos diferentes. Em redes viárias, um grafo G = (V, E)
representa vértices e arestas, eventualmente ponderadas por distância ou
tempo. A escolha de destinos pode utilizar regras sintéticas ou matrizes
origem--destino; porém, viagens agregadas entre regiões não oferecem,
isoladamente, descrição das trajetórias individuais. As hipóteses devem
ser declaradas quando os dados disponíveis possuem resolução espacial
mais ampla que o movimento simulado. (Taillandier; Drogoul, 2010;
Taillandier et al., 2019)

\subsubsection{2.5.2 Integração GIS--ABM: dados espaciais, geometrias e
agentificação}\label{integrauxe7uxe3o-gisabm-dados-espaciais-geometrias-e-agentificauxe7uxe3o}

A integração entre Sistemas de Informação Geográfica (SIG) e modelos
baseados em agentes permite que a representação do ambiente deixe de
depender exclusivamente de grades regulares ou coordenadas abstratas.
Taillandier e Drogoul (2010) distinguem três formas de incorporar dados
vetoriais à simulação: (i) leitura e escrita de arquivos geográficos;
(ii) utilização das feições como camadas espaciais que condicionam
deslocamentos e consultas; e (iii) agentificação, em que objetos
geográficos passam a ser instanciados como agentes com estados internos
e comportamentos. Essas modalidades não são equivalentes e devem ser
identificadas separadamente na descrição de cada modelo.

Na primeira modalidade, arquivos como shapefiles fornecem geometrias e
atributos utilizados na inicialização ou na exportação de resultados. Na
segunda, linhas, pontos e polígonos podem delimitar regiões acessíveis,
restringir movimentos ou fundamentar cálculos de proximidade e de
trajetórias. Na terceira, um objeto geográfico, como uma edificação ou
uma região, pode ser convertido em uma entidade computacional com
atributos mutáveis e regras próprias. A principal contribuição da
agentificação está na possibilidade de representar a evolução conjunta
da configuração espacial e do comportamento dos agentes (TAILLANDIER;
DROGOUL, 2010). (Taillandier; Drogoul, 2010; Taillandier et al., 2019)

A plataforma GAMA foi concebida para facilitar essas formas de
integração e oferece primitivas relacionadas a distâncias, geometrias,
vizinhanças e deslocamentos. O artigo apresenta essas capacidades em
aplicações de epidemiologia veterinária e resposta a desastres, nas
quais objetos espaciais participam do funcionamento da simulação.
Trata-se, contudo, de uma descrição histórica da plataforma: detalhes da
linguagem GAML, da interface e do conjunto de operadores devem ser
conferidos na documentação correspondente à versão efetivamente
utilizada no projeto. (Taillandier; Drogoul, 2010; Taillandier et al.,
2019)

Para o simulador de dengue deste trabalho, a incorporação de dados
geográficos do município de São Paulo é relevante para delimitar o
espaço da simulação e orientar a localização dos agentes. Já a
caracterização de uma eventual agentificação das unidades geográficas
depende de comprovação no código GAML: a simples leitura de um shapefile
ou a apresentação de polígonos no mapa não demonstram que cada feição
possui comportamentos autônomos. Também é necessário verificar projeção
cartográfica, unidades de distância e correspondência entre as regiões
representadas e os limites administrativos empregados na análise.
(Taillandier; Drogoul, 2010; Taillandier et al., 2019)

\subsection{2.6 Estocasticidade, incerteza e experimentos
reprodutíveis}\label{estocasticidade-incerteza-e-experimentos-reprodutuxedveis}

Em modelos estocásticos, eventos podem depender de números
pseudoaleatórios. Uma condição com probabilidade p pode ser avaliada
comparando um sorteio uniforme U(0,1) com p. Uma taxa contínua λ não
deve ser interpretada diretamente como probabilidade por ciclo: sob
hipótese de risco constante, a probabilidade de ocorrência no intervalo
Δt é p(Δt) = 1 − exp(−λΔt). Essa distinção é relevante quando o
intervalo simulado difere das unidades usadas nas fontes biológicas.
(Gamal et al., 2025; Morris, 1991)

A repetição de experimentos com diversas sementes caracteriza a
variabilidade aleatória. A reprodução de uma execução exige registrar
também versão do código, dados de entrada, parâmetros, plataforma,
cronograma de processos e ambiente computacional. Fixar a semente,
isoladamente, não resolve diferenças decorrentes da ordem de chamadas ao
gerador ou de alterações na execução. (Gamal et al., 2025; Morris, 1991)

É importante distinguir incerteza aleatória, associada aos sorteios do
próprio modelo, e incerteza epistêmica, associada a parâmetros e
mecanismos insuficientemente conhecidos. Análises de sensibilidade
investigam a influência das entradas; repetição simples não substitui
análise de adequação dos pressupostos. Morris (1991) propõe o uso de
efeitos elementares para triagem de parâmetros, enquanto Sobol (2001)
oferece uma abordagem de decomposição da variância, ambos condicionados
a um desenho experimental adequado.

\subsubsection{2.6.1 Experimentação computacional, parâmetros e
cenários}\label{experimentauxe7uxe3o-computacional-paruxe2metros-e-cenuxe1rios}

A experimentação computacional permite investigar como um modelo
responde à alteração de hipóteses, regras de interação, condições
iniciais e parâmetros. Em um modelo baseado em agentes, fenômenos
agregados podem emergir de ações locais, mas uma associação visual entre
determinada alteração e uma curva simulada não substitui um desenho
experimental controlado. É necessário definir antecipadamente fatores,
níveis, variáveis-resposta, número de replicações e critérios de
análise. (Morris, 1991; Gamal et al., 2025)

Sanderson (2020), na apresentação didática Simulating an epidemic,
exemplifica a comparação de cenários em um modelo SIR espacial
simplificado. O autor examina alterações no alcance das interações, nas
probabilidades de transmissão, na circulação entre comunidades e na
concentração de pessoas em pontos de encontro, observando diferenças na
evolução das curvas simuladas. A exposição também ilustra que execuções
com os mesmos parâmetros podem produzir trajetórias distintas devido à
aleatoriedade, reforçando a necessidade de replicações independentes.
Trata-se de demonstração pedagógica, e não de evidência de validade
epidemiológica do presente simulador.

Nesse contexto, convém distinguir três procedimentos: (i) experimentos
de cenário, que contrastam configurações previamente definidas; (ii)
análise de sensibilidade, que investiga quanto as respostas dependem de
entradas selecionadas; e (iii) avaliação estocástica, que estima a
dispersão das saídas quando variam as sementes pseudoaleatórias. A
alteração de um fator por vez facilita a interpretação inicial, mas não
identifica necessariamente efeitos de interação entre parâmetros. Para
investigações mais abrangentes, métodos como Morris (1991) e Sobol
(2001) oferecem referenciais de análise de sensibilidade, condicionados
à adequação do desenho amostral e do custo computacional.

A transposição desses exemplos para a dengue requer ressalvas. No modelo
SIR ilustrado no vídeo, a infecção é representada por interações diretas
entre pessoas; na dengue, a transmissão usual depende de mosquitos
vetores. Assim, parâmetros de distanciamento interpessoal, higiene
respiratória e isolamento não podem ser interpretados automaticamente
como medidas de controle vetorial. Para este trabalho, o aproveitamento
metodológico concentra-se na comparação controlada de parâmetros
pertinentes à interação humano--mosquito, na mobilidade espacial, na
distribuição dos vetores e na quantificação da variabilidade das
execuções. (Gamal et al., 2025; Morris, 1991)

\subsubsection{2.6.2 Fontes de incerteza em modelos baseados em
agentes}\label{fontes-de-incerteza-em-modelos-baseados-em-agentes}

A natureza estocástica de muitos modelos baseados em agentes introduz
variabilidade nas saídas mesmo quando as configurações de entrada são
mantidas constantes. Entretanto, a dispersão entre execuções não
representa a única fonte de incerteza. Gamal et al. (2025) distinguem,
no contexto de modelos comportamentais, incertezas associadas aos
parâmetros, à estrutura do modelo, à variabilidade entre execuções
estocásticas e às evidências empíricas utilizadas para representar o
sistema real.

A incerteza paramétrica decorre da seleção e dos valores atribuídos às
entradas; a incerteza estrutural relaciona-se às abstrações, regras e
interações adotadas na construção do modelo; a variabilidade entre
execuções resulta dos componentes aleatórios; e a incerteza empírica
está associada às limitações dos dados observados. Essas dimensões não
são equivalentes e não podem ser eliminadas apenas pela repetição das
simulações (GAMAL et al., 2025).

No estudo da transmissão da dengue, a distinção é relevante porque a
escolha das probabilidades de infecção, a representação de deslocamentos
e a resolução dos dados geográficos podem influenciar os resultados por
mecanismos diferentes. Por exemplo, executar o mesmo cenário com
múltiplas sementes permite caracterizar a variabilidade estocástica
condicionada à configuração adotada, mas não demonstra que os parâmetros
ou as regras implementadas correspondem aos processos epidemiológicos
reais. (Morris, 1991; Gamal et al., 2025)

A análise de sensibilidade (Sensitivity Analysis --- SA) e a
quantificação da incerteza (Uncertainty Quantification --- UQ)
constituem enfoques complementares. Enquanto a primeira investiga a
influência das entradas sobre as saídas, a segunda procura caracterizar
a incerteza dos resultados, conforme as fontes consideradas. A análise
local pode variar uma entrada de cada vez; abordagens globais exploram
combinações de parâmetros e permitem investigar interações que a
variação isolada pode deixar de captar (GAMAL et al., 2025).

Entre os métodos globais, os índices de Sobol permitem estimar
contribuições dos parâmetros para a variância das saídas, inclusive
efeitos associados às interações. Sua aplicação, contudo, requer um
desenho amostral apropriado e esforço computacional compatível com o
número de parâmetros e execuções. Portanto, a escolha do método deve
respeitar a questão de pesquisa, o orçamento computacional e a
estabilidade das estimativas, evitando interpretar alterações isoladas
de parâmetros como uma análise global completa. (Sobol, 2001; Gamal et
al., 2025)

\subsubsection{2.6.3 Heterogeneidade espacial e propagação temporal da
incerteza}\label{heterogeneidade-espacial-e-propagauxe7uxe3o-temporal-da-incerteza}

Em modelos espacialmente explícitos, os indicadores agregados podem
ocultar diferenças entre unidades geográficas. Duas execuções com totais
municipais semelhantes podem produzir distribuições regionais distintas.
Por isso, a avaliação da variabilidade apenas no encerramento das
simulações pode ser insuficiente para compreender onde e quando a
incerteza se manifesta. (Gamal et al., 2025; Morris, 1991)

Gamal et al. (2025) propõem a abordagem Spatio-Temporal Uncertainty
Analysis (ST-UA), que combina medidas de incerteza ou sensibilidade por
unidade espacial com índices destinados a caracterizar a relevância de
sua propagação entre dois instantes da simulação. Para demonstrar o
procedimento, os autores empregam o modelo Tobacco Town, voltado a
comportamentos relacionados ao tabagismo, e um indicador de
sensibilidade de efeitos totais associado ao método de Sobol. A
contribuição do artigo consiste na metodologia para examinar resultados
espacialmente heterogêneos, e não em regras ou parâmetros diretamente
transferíveis à transmissão da dengue.

A abordagem destaca a importância de definir antecipadamente a unidade
espacial, a variável de resposta e os instantes de comparação. No caso
de uma simulação aplicada ao município de São Paulo, análises futuras
poderiam considerar, mediante exportação adequada de dados, contagens ou
proporções de agentes infectados por unidades espaciais estáveis, em
diferentes tempos simulados. Tal aplicação dependeria da correspondência
entre as unidades geográficas usadas, da existência de replicações
suficientes e da consistência dos mecanismos de contagem. (Gamal et al.,
2025; Morris, 1991)

Uma estratégia metodológica progressiva consiste em começar pela
verificação das saídas e pela comparação descritiva entre execuções
repetidas, incluindo médias e medidas de dispersão ao longo do tempo.
Somente depois, se houver suporte nos dados e no desempenho da
implementação, seria justificável empregar índices globais de
sensibilidade ou a ST-UA completa. Mapas de incerteza, por si sós, não
validam epidemiologicamente o modelo; tornam explícita a variabilidade
das consequências de determinadas hipóteses computacionais. (Gamal et
al., 2025; Morris, 1991)

\subsection{2.7 Plataforma GAMA e linguagem
GAML}\label{plataforma-gama-e-linguagem-gaml}

A plataforma GAMA oferece recursos para especificar e executar modelos
espacialmente explícitos baseados em agentes. A linguagem GAML organiza
elementos como global, species, init, reflex, action e experiment,
permitindo definir estados, operações, inicialização e interfaces de
observação. Taillandier et al. (2019) descrevem sua aplicação à
construção, composição e experimentação de modelos espaciais complexos.
(Santos, A. R.; Santos, F., 2024)

A escolha do GAMA para este estudo decorre da necessidade de articular
agentes heterogêneos, dados geográficos e observação das execuções. Essa
adequação funcional não constitui demonstração de superioridade de
desempenho sobre NetLogo, Mesa, Repast ou AnyLogic. Comparações de
velocidade ou memória exigiriam implementações e condições experimentais
equivalentes, não examinadas nesta pesquisa. (Taillandier et al., 2019;
Taillandier; Drogoul, 2010)

\subsubsection{2.7.1 Primitivas geométricas e relevância para interações
espaciais}\label{primitivas-geomuxe9tricas-e-relevuxe2ncia-para-interauxe7uxf5es-espaciais}

Taillandier e Drogoul (2010) descrevem operações geométricas
disponibilizadas no GAMA, entre elas cálculo de distância, teste de
interseção, construção de áreas de influência, seleção de pontos dentro
de polígonos, identificação de vizinhos e cálculo de trajetórias em
geometrias ou redes. Essas possibilidades são especialmente pertinentes
à representação de encontros entre agentes humanos e mosquitos em
modelos de transmissão vetorial. Entretanto, a disponibilidade de um
operador na plataforma não autoriza afirmar que ele foi utilizado no
simulador analisado; tal afirmação exige rastreabilidade até o
código-fonte.

Além disso, consultas espaciais envolvem escolhas de implementação com
consequências para o desempenho. A interpretação de um raio de busca
exige identificar a unidade do sistema de referência espacial e a
semântica do operador de distância; uma distância numérica sem unidade
explicitada não representa automaticamente metros no território real. A
eventual utilização de estruturas de indexação, grafos ou algoritmos de
caminho mínimo deve ser apresentada somente após inspeção da
implementação correspondente. (Taillandier et al., 2019; Taillandier;
Drogoul, 2010)

\subsection{2.8 Documentação ODD e rastreabilidade entre modelo e
software}\label{documentauxe7uxe3o-odd-e-rastreabilidade-entre-modelo-e-software}

O protocolo ODD (Overview, Design concepts and Details) organiza a
descrição de modelos baseados em agentes em componentes relacionados a
propósito, entidades, escalas, processos, inicialização, entradas e
submodelos. Grimm et al. (2020) ressaltam a importância de tornar
descrições compreensíveis, verificáveis e úteis à replicação, incluindo
as justificativas das escolhas de modelagem.

Neste TCC, o ODD será utilizado como referência para mapear requisitos
conceituais aos componentes GAML, explicando o que o software realmente
faz. O documento da especificação, entretanto, não substitui a análise
do código, nem comprova sozinho a correção dos processos. A
rastreabilidade depende de trechos identificáveis, versões documentadas
e testes executáveis. (Grimm et al., 2020)

\subsection{2.9 Verificação, validação e desempenho
computacional}\label{verificauxe7uxe3o-validauxe7uxe3o-e-desempenho-computacional}

O rigor de verificação e validação de modelos epidemiológicos baseados
em agentes é discutido por Arifin, Madey e Collins (2016), que
distinguem atividades destinadas a identificar erros de implementação
daquelas que buscam avaliar a adequação do modelo à realidade
investigada. Essas dimensões também se distinguem da reprodutibilidade:
resultados comparáveis requerem o registro da versão do software,
parâmetros, sementes e ambiente de execução. No presente estudo, a
verificação interna constitui requisito anterior a qualquer alegação de
validação externa.

Verificação examina se a implementação corresponde às regras
especificadas; validação avalia em que medida o modelo é adequado ao uso
pretendido em relação ao sistema representado; calibração corresponde ao
ajuste de parâmetros perante observações. Os procedimentos precisam ser
distintos: um programa pode funcionar sem erros aparentes e produzir
padrões epidemiológicos inadequados. Da mesma maneira, a proximidade
entre uma curva simulada e outra observada não confirma isoladamente a
correção dos mecanismos internos. (Grimm et al., 2020; Macal; North,
2010)

A avaliação computacional pode utilizar cenários controlados para
verificar estados e invariantes, testes de reprodutibilidade, comparação
entre sementes e perfis de execução. Indicadores de desempenho incluem
tempo de inicialização, duração dos ciclos, memória e custo do registro
de saídas. No contexto de São Paulo, importa separar escala geográfica
da quantidade de agentes sintéticos, evitando interpretar redução
populacional como equivalência automática ao comportamento do município
real. (Grimm et al., 2020; Macal; North, 2010)

\subsection{2.10 Trabalhos relacionados}\label{trabalhos-relacionados}

\subsubsection{2.10.1 Garcias (2024): simulador GAMA em
Portugal}\label{garcias-2024-simulador-gama-em-portugal}

Garcias (2024) apresenta um modelo baseado em agentes para a transmissão
da dengue no contexto de Portugal Continental, utilizando principalmente
Aedes albopictus, agentes humanos e entidades territoriais. O trabalho
descreve decisões epidemiológicas, parâmetros, mobilidade e componentes
ambientais, além de organizar cenários e procedimentos de verificação e
validação. Sua principal contribuição para este TCC é a estrutura de
documentação dos comportamentos e a relação entre representação
geográfica, regras e experimentação.

A diferença entre os vetores e as áreas de aplicação impede a
transferência automática de parâmetros, dados demográficos e conclusões.
Além disso, a afirmação de validação do estudo anterior não equivale à
validação do software analisado neste trabalho. (Garcias, 2024; Oliveira
et al., 2025)

\subsubsection{2.10.2 Oliveira et al. (2025): simulação multiagentes em
Alto
Santo}\label{oliveira-et-al.-2025-simulauxe7uxe3o-multiagentes-em-alto-santo}

Oliveira et al. (2025) apresentam uma simulação multiagentes da dengue
aplicada ao município de Alto Santo, no Ceará, também na plataforma
GAMA. Os autores representam humanos, mosquitos, ovos e criadouros;
empregam dados geoespaciais e realizam experimentos variando mortalidade
vetorial, oviposição, criadouros, estratégias de nebulização e
distribuição espacial. Os resultados apresentados pertencem ao contexto
daquele estudo e não devem ser tratados como resultados da simulação de
São Paulo.

O trabalho fornece um exemplo concreto de organização de parâmetros e de
comparação de cenários. Seu horizonte experimental e suas hipóteses
diferem do presente projeto; por isso, sua utilidade está no desenho da
investigação e na comparação crítica, não na substituição do protocolo
de verificação próprio. (Garcias, 2024; Oliveira et al., 2025)

\subsubsection{2.10.3 Modelos individuais de dengue e de outras doenças
transmitidas por
mosquitos}\label{modelos-individuais-de-dengue-e-de-outras-doenuxe7as-transmitidas-por-mosquitos}

Isidoro et al. (2011) desenvolveram um modelo baseado em agentes voltado
à interação entre populações humanas e populações de Aedes aegypti na
transmissão da dengue. O estudo é particularmente pertinente ao presente
TCC por explicitar a representação dos vetores e das relações de
transmissão em uma população simulada. Entretanto, a plataforma, a
estrutura do modelo e os parâmetros daquele trabalho não equivalem aos
adotados no simulador GAMA investigado nesta pesquisa; seu emprego aqui
é metodológico e comparativo, e não uma fonte de parâmetros
automaticamente transferíveis.

Jindal e Rao (2017) propuseram um modelo de transmissão de doenças
veiculadas por mosquitos que articula a mobilidade dos hospedeiros, a
distribuição espacial das populações, os dados geográficos e as
condições climáticas. Os autores apresentam uma aplicação relacionada à
chikungunya e discutem a utilização de dados epidemiológicos para
avaliação do modelo. Sua contribuição reforça o interesse de combinar
representação individual e ambiente espacial; todavia, qualquer
capacidade preditiva obtida naquele estudo não implica que o simulador
de dengue desenvolvido neste TCC seja preditivo ou tenha sido validado
empiricamente.

Em termos comparativos, os dois estudos ajudam a distinguir três
decisões de projeto: como representar entidades humanas e vetoriais;
quais dados espaciais e ambientais alimentam as regras; e como testar se
os resultados decorrem de uma implementação correta. Tais decisões
orientam a especificação e a verificação do modelo deste trabalho, mas
somente podem ser consideradas implementadas quando confirmadas no
código-fonte e nos registros experimentais correspondentes (Isidoro et
al., 2011; Jindal; Rao, 2017).

\subsubsection{2.10.4 Estudos espaciais de vetores e outras
doenças}\label{estudos-espaciais-de-vetores-e-outras-doenuxe7as}

Arifin, Madey e Collins (2016) apresentam uma abordagem sistemática de
simulação espacial baseada em agentes aplicada à epidemiologia da
malária. A obra articula processos biológicos, agentes vetoriais e
informações geográficas, incluindo discussões específicas de
verificação, validação, replicação e reprodutibilidade. Embora trate
sobretudo de Anopheles e malária, e não de Aedes aegypti e dengue, sua
discussão de engenharia de modelos e avaliação computacional é
aproveitável no presente projeto, sem transpor os parâmetros biológicos
entre doenças.

Gharakhanlou, Mesgari e Hooshangi (2019) e Gharakhanlou, Hooshangi e
Helbich (2020) discutem modelos espaciais baseados em agentes para
malária na região de Sarbaz, no Irã. Embora o sistema biológico seja
diferente, os trabalhos contribuem para compreender a integração entre
entidades, espaço, parâmetros ambientais e intervenções. As regras
epidemiológicas específicas da malária não são transferíveis sem
adaptação ao mecanismo da dengue. (Gharakhanlou; Mesgari; Hooshangi,
2019; Gharakhanlou; Hooshangi; Helbich, 2020)

\subsubsection{2.10.5 Gamal et al. (2025): propagação espaço-temporal da
incerteza}\label{gamal-et-al.-2025-propagauxe7uxe3o-espauxe7o-temporal-da-incerteza}

Gamal et al. (2025) apresentam uma abordagem de análise da incerteza
espaço-temporal aplicada a modelos baseados em agentes. Diferentemente
de estudos centrados exclusivamente na evolução de um indicador
agregado, sua proposta examina como medidas de sensibilidade ou
incerteza variam entre unidades espaciais e entre estados temporais do
modelo.

O estudo emprega o Tobacco Town, um modelo de comportamentos
relacionados ao tabagismo, para ilustrar a metodologia ST-UA. A
aplicação utiliza distribuições de variáveis individuais, múltiplas
execuções e índices relacionados à decomposição da variância, com
resultados representados espacialmente. A originalidade metodológica
reside na comunicação e identificação de áreas com alterações relevantes
na incerteza, não em uma contribuição específica à epidemiologia da
dengue. (Gamal et al., 2025)

Para o presente TCC, o artigo oferece fundamentação para tratar a
variabilidade estocástica e a heterogeneidade espacial como dimensões
distintas da avaliação computacional. Sua metodologia poderá inspirar
protocolos de observação por região e tempo, mas a reprodução integral
dos índices propostos não constitui resultado alcançado e exigiria
adequação do armazenamento de dados, definição das unidades de análise,
controle de parâmetros e recursos computacionais. (Gamal et al., 2025)

\subsubsection{2.10.6 Comparação crítica e posicionamento do presente
trabalho}\label{comparauxe7uxe3o-cruxedtica-e-posicionamento-do-presente-trabalho}

Como trabalho brasileiro recente, Doege (2026) apresenta um TCC de
Biomedicina na UFCSPA que utilizou NetLogo, 500 agentes humanos, 250
mosquitos e estados SEIR/SEI para comparar estratégias de controle da
dengue. O autor documenta cinco cenários de intervenção e 30 execuções
independentes por cenário. Trata-se de uma referência metodológica para
definir métricas de pico, duração e carga acumulada, além de organizar
réplicas estocásticas; suas conclusões e seus valores quantitativos não
podem ser atribuídos às simulações do presente TCC, que utilizam outro
software e requerem testes próprios.

De forma convergente, Isidoro et al. (2011) enfatizam a representação do
Aedes aegypti, Jindal e Rao (2017) enfatizam a integração de dados
espaciais, ambientais e de mobilidade, enquanto Arifin, Madey e Collins
(2016) sistematizam mecanismos de avaliação e reprodutibilidade. A
contribuição diferenciadora buscada aqui é documentar e examinar
tecnicamente um simulador de dengue na plataforma GAMA aplicado ao
município de São Paulo, sem reivindicar inovação epidemiológica ou
capacidade preditiva ainda não demonstradas.

Quadro 3 --- Síntese comparativa dos trabalhos relacionados

\begin{longtable}[]{@{}
  >{\raggedright\arraybackslash}p{(\linewidth - 6\tabcolsep) * \real{0.2500}}
  >{\raggedright\arraybackslash}p{(\linewidth - 6\tabcolsep) * \real{0.2500}}
  >{\raggedright\arraybackslash}p{(\linewidth - 6\tabcolsep) * \real{0.2500}}
  >{\raggedright\arraybackslash}p{(\linewidth - 6\tabcolsep) * \real{0.2500}}@{}}
\toprule\noalign{}
\begin{minipage}[b]{\linewidth}\raggedright
\textbf{Estudo}
\end{minipage} & \begin{minipage}[b]{\linewidth}\raggedright
\textbf{Domínio e enfoque}
\end{minipage} & \begin{minipage}[b]{\linewidth}\raggedright
\textbf{Contribuição ao presente TCC}
\end{minipage} & \begin{minipage}[b]{\linewidth}\raggedright
\textbf{Não transferir automaticamente}
\end{minipage} \\
\midrule\noalign{}
\endhead
\bottomrule\noalign{}
\endlastfoot
Garcias (2024) & Dengue / A. albopictus / Portugal / GAMA &
Especificação dos agentes, cenários e documentação & Biologia e
parâmetros do vetor \\
Oliveira et al. (2025) & Dengue / A. aegypti / Alto Santo / GAMA &
Desenho experimental e análises de parâmetros & Resultados e
configuração territorial \\
Gharakhanlou et al. (2019; 2020) & Malária / Sarbaz / ABM espacial &
Integração geográfica e discussão de intervenções & Regras de
transmissão da malária \\
Este TCC & Dengue / São Paulo / GAMA & Análise de código,
rastreabilidade, verificação e avaliação computacional & A validade
ainda depende de evidências próprias \\
\end{longtable}

Fonte: elaboração própria, a partir das referências citadas.

A contribuição que este TCC pode efetivamente demonstrar não decorre
apenas da mudança do território para São Paulo. Ela deverá resultar da
descrição rastreável da implementação, da identificação de escolhas e
limitações computacionais, da execução de testes e da apresentação de
resultados próprios. O repositório público constitui fonte primária para
a descrição técnica posterior, mas as afirmações de correção e
desempenho dependem de evidência adicional. (Taillandier et al., 2019;
Gamal et al., 2025)

\subsubsection{2.10.7 Taillandier e Drogoul (2010): integração de GIS e
agentificação na plataforma
GAMA}\label{taillandier-e-drogoul-2010-integrauxe7uxe3o-de-gis-e-agentificauxe7uxe3o-na-plataforma-gama}

Taillandier e Drogoul (2010) apresentam uma contribuição conceitual para
a integração de informações geográficas em modelos baseados em agentes,
enfatizando a possibilidade de representar feições vetoriais como
agentes. Os exemplos de propagação de influenza aviária e de resposta a
desastres demonstram a aplicabilidade dessa abordagem em sistemas cujas
relações dependem do espaço. Em contraste com trabalhos voltados à
dinâmica epidemiológica específica da dengue, essa referência fundamenta
a arquitetura espacial e os recursos da plataforma. Sua principal
limitação para este TCC é a antiguidade tecnológica: a comparação com
plataformas concorrentes e as descrições de versões do GAMA refletem o
contexto de publicação de 2010 e não substituem uma avaliação atual das
ferramentas.

\subsection{2.11 Síntese do capítulo}\label{suxedntese-do-capuxedtulo}

A literatura examinada fundamenta a escolha da modelagem baseada em
agentes para representar heterogeneidade individual e interações
espaciais na transmissão da dengue. Também estabelece requisitos
computacionais: definição precisa das máquinas de estados, integração de
dados territoriais, controle temporal, rastreabilidade dos algoritmos e
reprodução de experimentos. O capítulo seguinte transforma esses
conceitos em uma especificação explícita do modelo, distinguindo o
comportamento desejado das regras presentes na implementação examinada.
(Taillandier et al., 2019; Gamal et al., 2025)

A discussão de Gamal et al. (2025) complementa essa fundamentação ao
demonstrar que a avaliação de modelos espaciais pode exigir a
caracterização das incertezas tanto entre regiões quanto ao longo do
tempo. No presente estudo, tal perspectiva constitui orientação
metodológica para análises futuras, e não evidência de experimentos já
executados.

\section{3 MODELAGEM COMPUTACIONAL E JUSTIFICATIVA DA
PLATAFORMA}\label{modelagem-computacional-e-justificativa-da-plataforma}

Este capítulo estabelece a especificação computacional do simulador de
transmissão da dengue aplicado ao município de São Paulo. O propósito
não é repetir a fundamentação epidemiológica, mas transformá-la em
requisitos, entidades, estados, regras, entradas e saídas verificáveis.
Também apresenta os critérios utilizados para justificar a seleção da
plataforma GAMA e dos recursos de Sistemas de Informação Geográfica (SIG
ou GIS), confrontando-os com alternativas reconhecidas de modelagem
baseada em agentes. (Taillandier; Drogoul, 2010; Taillandier et al.,
2019)

A especificação distingue três níveis de evidência: regras desejadas do
modelo conceitual; comportamentos reconhecidos nos arquivos públicos do
projeto TccAbm; e procedimentos de avaliação ainda necessários. A
presença de uma função no código não comprova sua execução correta, e
saídas armazenadas não constituem, isoladamente, validação
epidemiológica. (HKIOSHI, 2026)

\subsection{3.1 Propósito, fronteira e requisitos da
simulação}\label{propuxf3sito-fronteira-e-requisitos-da-simulauxe7uxe3o}

O sistema investigado é uma simulação espacial e estocástica da
interação entre pessoas, mosquitos Aedes aegypti e estruturas
ambientais. O ambiente territorial corresponde ao município de São
Paulo. A representação não pretende reproduzir um a um os habitantes
reais: agentes sintéticos ocupam regiões do município conforme regras de
inicialização e dados geográficos. Consequentemente, devem ser
distinguidos o domínio espacial, a resolução das unidades cartográficas
e o tamanho efetivo da população representada. (Garcias, 2024; Oliveira
et al., 2025)

O objetivo computacional é permitir observação, explicação e comparação
de trajetórias simuladas sob alterações de parâmetros. A utilização
epidemiológica é exploratória: padrões produzidos dependem de hipóteses,
escalas, condições iniciais e algoritmos, não devendo ser interpretados
como previsões de notificações oficiais sem calibração e validação
independentes. (Macal; North, 2010; Grimm et al., 2020)

Quadro 4 --- Requisitos computacionais e evidências esperadas

\begin{longtable}[]{@{}
  >{\raggedright\arraybackslash}p{(\linewidth - 4\tabcolsep) * \real{0.3333}}
  >{\raggedright\arraybackslash}p{(\linewidth - 4\tabcolsep) * \real{0.3333}}
  >{\raggedright\arraybackslash}p{(\linewidth - 4\tabcolsep) * \real{0.3333}}@{}}
\toprule\noalign{}
\begin{minipage}[b]{\linewidth}\raggedright
\textbf{Código}
\end{minipage} & \begin{minipage}[b]{\linewidth}\raggedright
\textbf{Requisito}
\end{minipage} & \begin{minipage}[b]{\linewidth}\raggedright
\textbf{Evidência ou teste}
\end{minipage} \\
\midrule\noalign{}
\endhead
\bottomrule\noalign{}
\endlastfoot
R01 & Criar e identificar diferentes espécies de agentes & Inspeção das
species e testes de contagem \\
R02 & Carregar unidades geográficas compatíveis com o município &
Inventário de arquivos GIS, projeção e cobertura \\
R03 & Definir transições epidemiológicas sem estados inválidos & Tabela
de transições e testes de invariantes \\
R04 & Implementar contatos espaciais e deslocamentos & Trajetórias e
inspeção de distâncias \\
R05 & Controlar tempo simulado e probabilidades & Rastreio dos ciclos e
teste de frequência \\
R06 & Registrar parâmetros, sementes e indicadores & Arquivos de saída e
metadados da execução \\
R07 & Viabilizar execuções comparáveis e aferir custo & Protocolos de
repetição, tempo e memória \\
\end{longtable}

Fonte: elaboração própria.

\subsection{3.2 Seleção do ambiente de desenvolvimento: GAMA com
GIS}\label{seleuxe7uxe3o-do-ambiente-de-desenvolvimento-gama-com-gis}

\subsubsection{3.2.1 Critérios de
seleção}\label{crituxe9rios-de-seleuxe7uxe3o}

A definição da ferramenta foi orientada pela compatibilidade com as
características do problema: heterogeneidade de agentes; geometrias
reais; proximidade e mobilidade; regras temporais e probabilísticas;
visualização do estado do sistema; experimentos parametrizados;
disponibilidade para uso acadêmico; e facilidade de registrar a
implementação. Esses critérios dizem respeito à adequação funcional, e
não à superioridade universal de uma linguagem ou plataforma.
(Taillandier; Drogoul, 2010; Taillandier et al., 2019) (Santos, A. R.;
Santos, F., 2024)

GAMA disponibiliza uma linguagem específica de modelagem, GAML, na qual
elementos como global, species, init, reflex, action e experiment
oferecem uma organização diretamente relacionada aos conceitos de
modelos baseados em agentes. A plataforma possui orientação geoespacial
e recursos para representar objetos e consultas no mesmo ambiente da
simulação (TAILLANDIER et al., 2019). A documentação de experimentos em
lote descreve mecanismos para explorar configurações de parâmetros e
replicações com diferentes sementes (GAMA PLATFORM, s.d.).

\subsubsection{3.2.2 Alternativas
consideradas}\label{alternativas-consideradas}

Banos, Lang e Marilleau (2015) apresentam um guia metodológico de
simulação espacial baseada em agentes utilizando NetLogo, com ênfase na
concepção de modelos, exploração de cenários e boas práticas de
modelagem. A referência contribui para contextualizar NetLogo como
alternativa relevante no campo, mas não constitui evidência de
superioridade ou inferioridade de desempenho da plataforma GAMA,
comparação que exigiria benchmarks sob condições equivalentes.

NetLogo oferece um ambiente amplamente difundido para experimentação com
agentes. Sua extensão GIS carrega dados vetoriais, incluindo Shapefile e
GeoJSON, e grades raster, mediante transformação entre coordenadas
geográficas e espaço do modelo. Portanto, seria incorreto dizer que
NetLogo não trabalha com dados geográficos; a diferença pertinente é a
forma de organizar a integração espacial e o fluxo de desenvolvimento.
(Taillandier; Drogoul, 2010; Taillandier et al., 2019)

Mesa, em Python, permite estruturar agentes em um ecossistema de
programação geral, enquanto Mesa-Geo oferece estruturas geográficas e
integração com bibliotecas de análise espacial. Repast Simphony fornece
recursos baseados em Java e objetos geográficos; AnyLogic permite
combinar paradigmas de simulação e possui recursos de mapas e
deslocamentos, sujeitos às funcionalidades e licenciamento da edição
escolhida. Todas constituem alternativas tecnicamente plausíveis, mas
exigiriam decisões distintas de implementação, dependências e
documentação (NETLOGO, 2026; MESA, s.d.; MESA-GEO, s.d.; REPAST, s.d.;
ANYLOGIC, s.d.). (NetLogo, {[}s. d.{]}; Mesa, {[}s. d.{]}; Repast, {[}s.
d.{]}; AnyLogic, {[}s. d.{]})

Quadro 5 --- Comparação qualitativa de plataformas para o projeto

\begin{longtable}[]{@{}
  >{\raggedright\arraybackslash}p{(\linewidth - 6\tabcolsep) * \real{0.2500}}
  >{\raggedright\arraybackslash}p{(\linewidth - 6\tabcolsep) * \real{0.2500}}
  >{\raggedright\arraybackslash}p{(\linewidth - 6\tabcolsep) * \real{0.2500}}
  >{\raggedright\arraybackslash}p{(\linewidth - 6\tabcolsep) * \real{0.2500}}@{}}
\toprule\noalign{}
\begin{minipage}[b]{\linewidth}\raggedright
\textbf{Ferramenta}
\end{minipage} & \begin{minipage}[b]{\linewidth}\raggedright
\textbf{GIS e representação espacial}
\end{minipage} & \begin{minipage}[b]{\linewidth}\raggedright
\textbf{Programação e experimentos}
\end{minipage} & \begin{minipage}[b]{\linewidth}\raggedright
\textbf{Implicação para este estudo}
\end{minipage} \\
\midrule\noalign{}
\endhead
\bottomrule\noalign{}
\endlastfoot
GAMA & Integração geográfica no ambiente de simulação & GAML;
experimentos gráficos e em lote & Alinhamento direto com o modelo
geográfico já implementado \\
NetLogo & Extensão GIS, Shapefile, GeoJSON e raster & Linguagem NetLogo;
BehaviorSpace & Alternativa didática; requer adequação da integração
espacial \\
Mesa + Mesa-Geo & Agentes geográficos no ecossistema GeoPandas & Python
e bibliotecas auxiliares & Flexibilidade e análise; mais decisões de
integração \\
Repast Simphony & Representação geográfica no ecossistema Java & Java e
framework de simulação & Alternativa de engenharia com curva própria \\
AnyLogic & GIS e roteamento conforme recursos disponíveis & Modelagem
gráfica e Java & Opção híbrida; avaliar limites de licença/edição \\
\end{longtable}

Fonte: elaboração própria, com base nas documentações das plataformas;
comparação funcional, não benchmark.

\subsubsection{3.2.3 Justificativa da adoção e
limitações}\label{justificativa-da-adouxe7uxe3o-e-limitauxe7uxf5es}

A adoção de GAMA com GIS mostrou-se coerente com a necessidade de reunir
no mesmo modelo agentes humanos, mosquitos e entidades territoriais,
cuja localização deve participar das operações de inicialização,
movimento e contato. O repositório analisado apresenta arquivos GAML
separados e um ponto de entrada Main.gaml, além de carregar camadas
geográficas, constituindo evidência prática de que essa tecnologia
corresponde à arquitetura escolhida no projeto (HKIOSHI, s.d.).
(HKIOSHI, 2026)

Os dados GIS não são um programa independente que ``executa'' a
simulação. Eles fornecem geometrias, atributos e relações espaciais;
GAMA usa essas informações para instanciar entidades e calcular
interações. O funcionamento depende da qualidade das camadas, da
adequação do sistema de coordenadas, dos parâmetros de distância e do
custo de processamento. Sem implementação equivalente nas alternativas,
não se conclui que GAMA seja mais veloz ou precise de menos memória.
(Taillandier; Drogoul, 2010; Taillandier et al., 2019)

\subsubsection{3.2.4 Integração de dados GIS e níveis de representação
espacial}\label{integrauxe7uxe3o-de-dados-gis-e-nuxedveis-de-representauxe7uxe3o-espacial}

A escolha do GAMA também se fundamenta na integração de dados vetoriais
com entidades e operações de simulação. Taillandier e Drogoul (2010)
distinguem a importação de informações geográficas, o emprego de feições
como suporte espacial e a agentificação das feições, pela qual cada
objeto geográfico pode adquirir estados e comportamentos próprios. Tal
distinção é importante para evitar a interpretação de qualquer camada
vetorial carregada como uma população de agentes autônomos.

Na aplicação ao município de São Paulo, a leitura de um shapefile e a
organização espacial dos agentes constituem funcionalidades cuja
implementação pode ser inspecionada no código. A classificação das
unidades territoriais como agentes, por sua vez, exige verificar a
criação de espécies e instâncias correspondentes e suas regras de
execução. A utilização de regiões administrativas, geometria de
referência e densidades sintéticas deve ser explicitada como escolha de
modelagem, não como reprodução integral da população ou da transmissão
real. (HKIOSHI, 2026)

A adequação espacial do simulador depende, adicionalmente, da
consistência do sistema de coordenadas e das unidades utilizadas nas
operações de proximidade. Por isso, o protocolo de verificação deverá
conferir a correspondência entre os dados geográficos de entrada, a
escala do mapa, as distâncias efetivamente calculadas e os limites do
município. Operações de caminho mínimo e algoritmos de grafos citados na
literatura sobre GAMA são possibilidades da plataforma; sua presença no
modelo permanece condicionada a evidências no código-fonte. (HKIOSHI,
2026)

\subsection{3.3 Delimitação territorial e representação
espacial}\label{delimitauxe7uxe3o-territorial-e-representauxe7uxe3o-espacial}

A área de estudo é o município de São Paulo. A divisão em Centro, Norte,
Leste, Sul e Oeste aparece nas regras de inicialização previamente
identificadas no código, enquanto o arquivo denominado
subprefeitura\_v2.shp fornece unidades geográficas ao modelo. Essas
designações não devem ser confundidas sem conferir suas geometrias,
extensão territorial, atributos e regras de associação. A existência de
um Shapefile de subprefeituras não basta para demonstrar que a
representação cobre adequadamente todos os distritos ou que as
proporções populacionais estão atualizadas. (HKIOSHI, 2026)

A preparação da base requer documentar sua fonte, ano, sistema de
referência de coordenadas (SRC), unidades de comprimento, consistência
topológica, critérios de agregação territorial e tratamento de
geometrias problemáticas. Caso as operações de distância utilizem
coordenadas angulares em vez de métricas, um raio expresso no código não
poderá ser interpretado diretamente em metros. Os critérios serão
verificados na etapa de implementação e no protocolo experimental.
(HKIOSHI, 2026)

\subsection{3.4 Entidades, atributos e
relações}\label{entidades-atributos-e-relauxe7uxf5es}

Os agentes humanos representam hospedeiros com localização, destino,
estado epidemiológico, condição vacinal e rotinas. Os mosquitos
representam vetores que podem movimentar-se, procurar humanos e
participar de eventos probabilísticos de infecção. Elementos ambientais
e criadouros fornecem suporte à dinâmica da população vetorial. O nível
de individualização dos estágios imaturos deve ser descrito a partir do
conteúdo efetivo dos módulos do repositório, sem assumir a existência de
uma species para cada fase biológica. (HKIOSHI, 2026)

Quadro 6 --- Entidades e variáveis para a especificação

\begin{longtable}[]{@{}
  >{\raggedright\arraybackslash}p{(\linewidth - 4\tabcolsep) * \real{0.3333}}
  >{\raggedright\arraybackslash}p{(\linewidth - 4\tabcolsep) * \real{0.3333}}
  >{\raggedright\arraybackslash}p{(\linewidth - 4\tabcolsep) * \real{0.3333}}@{}}
\toprule\noalign{}
\begin{minipage}[b]{\linewidth}\raggedright
\textbf{Entidade}
\end{minipage} & \begin{minipage}[b]{\linewidth}\raggedright
\textbf{Principais atributos / variáveis}
\end{minipage} & \begin{minipage}[b]{\linewidth}\raggedright
\textbf{Operações ou relações}
\end{minipage} \\
\midrule\noalign{}
\endhead
\bottomrule\noalign{}
\endlastfoot
Humano & Identificador; posição; residência e destino; estado; vacinação
& Movimento, contato, progressão de estados \\
Mosquito & Posição; idade; estado; parâmetros de picada & Movimento,
picada, infecção, mortalidade \\
Criador/poça & Geometria ou posição; água; capacidade & Atualização
ambiental e suporte à reprodução \\
Unidade territorial & Geometria; identificador; região; atributos &
Pertencimento, distribuição inicial e consultas \\
Controlador global & Relógio; parâmetros; semente; contadores &
Inicialização, eventos, exportação \\
\end{longtable}

Fonte: elaboração própria.

A expressão ``criador/poça'' neste quadro é uma classe conceitual: o
identificador exato e a existência de objetos específicos devem ser
confirmados no arquivo Poca.gaml. De modo semelhante, os atributos acima
não são uma transcrição literal completa das variáveis; constituem o
conjunto necessário à documentação funcional. (HKIOSHI, 2026)

\subsection{3.5 Estados epidemiológicos e máquinas de
estados}\label{estados-epidemioluxf3gicos-e-muxe1quinas-de-estados}

A especificação de referência considera humanos com estados suscetível
(S), exposto (E), infectante (I) e recuperado (R), e mosquitos adultos
com estados suscetível (S), exposto (E) e infectante (I). A condição
vacinal é independente do estado epidemiológico. As transições precisam
corresponder a eventos explícitos; um contato isolado não implica
transmissão, e a incubação não deve ser confundida com infectividade.
(Garcias, 2024; Oliveira et al., 2025)

Humano: S → E → I → R

Mosquito: S → E → I (com possibilidade de morte em qualquer estágio)

Essa formulação representa o modelo conceitual adotado para análise, não
uma certificação de correspondência integral com o código. A inspeção
anterior identificou uso de rnd(0,3) na inicialização vetorial e uma
operação de transição circular (estado + 1) mod 3. É necessário testar
se essas rotinas introduzem estados não previstos ou permitem retorno
inadequado do mosquito infectante ao estado suscetível. A seção de
implementação deverá registrar literalmente as condições de cada
transição, os períodos e os gatilhos temporais. (HKIOSHI, 2026)

\subsection{3.6 Interações, transmissão e
probabilidade}\label{interauxe7uxf5es-transmissuxe3o-e-probabilidade}

Uma oportunidade de transmissão exige um mosquito e um humano elegíveis,
proximidade espacial, ocorrência da picada e sucesso probabilístico de
transmissão conforme a direção humano--mosquito ou mosquito--humano.
Esses eventos devem ser especificados separadamente para distinguir
frequência de contatos e probabilidade condicional de infecção. O
tratamento da condição vacinal precisa indicar a variável de proteção e
seu alvo: infecção, manifestação clínica ou gravidade. (Garcias, 2024;
Oliveira et al., 2025)

P(infecção \textbar{} contato elegível) = p\_base × (1 − ε\_inf)
{[}hipótese de redução por vacinação{]}

A equação anterior é uma regra conceitual para proteção contra infecção
e somente deve ser atribuída à implementação se sua lógica e parâmetros
forem confirmados no código. Taxas definidas por dia não podem ser
reaplicadas sem conversão em múltiplos ciclos subdiários; sob hipótese
de risco exponencial constante, p(Δt) = 1 − exp(−λΔt). A adequação dessa
hipótese depende do processo e do intervalo temporal representados.
(HKIOSHI, 2026)

\subsection{3.7 Mobilidade humana e consultas de
vizinhança}\label{mobilidade-humana-e-consultas-de-vizinhanuxe7a}

A modelagem da mobilidade diferencia decisão de destino e percurso
realizado. As rotinas reconhecidas anteriormente em Pessoa.gaml
contemplam janelas de saída de 5h a 8h e retorno de 16h a 19h, além de
deslocamentos locais. A documentação definitiva deverá esclarecer
seleção do destino, velocidades, uso de redes, restrições geográficas e
o significado da unidade temporal. Atribuir viagens sem verificar a
conectividade das vias pode gerar trajetórias irreais. (HKIOSHI, 2026)

As consultas de proximidade utilizadas em eventos de picada são
particularmente sensíveis à resolução espacial e ao número de agentes.
Se uma busca examinar indiscriminadamente todos os pares, seu custo
potencial cresce como O(N\_H × N\_M). A existência de operações
espaciais em GAML não dispensa medir seu custo real, que dependerá dos
filtros usados, da distribuição dos agentes e das estruturas de
indexação disponíveis. (HKIOSHI, 2026)

\subsection{3.8 Ambiente, clima, desenvolvimento e reprodução
vetorial}\label{ambiente-clima-desenvolvimento-e-reproduuxe7uxe3o-vetorial}

O módulo Poca.gaml contém regras de atualização de poças em função de
condições ambientais como chuva e temperatura, segundo a análise
anterior do repositório. Para que essas relações sejam reprodutíveis,
será preciso identificar parâmetros, unidades, frequência de
atualização, fonte meteorológica e interação com os mecanismos de
reprodução. Uma relação climática implementada como regra simplificada
não constitui estimativa observacional direta de larvas ou mosquitos.
(HKIOSHI, 2026)

A representação dos estágios imaturos exige verificar se ovos, larvas e
pupas são agentes, grupos ou estados de uma mesma entidade. É
metodologicamente preferível documentar a estrutura existente e
explicitar simplificações do que descrever um ciclo detalhado que ainda
não corresponda aos objetos GAML. (HKIOSHI, 2026)

\subsection{3.9 Tempo, agendamento e
aleatoriedade}\label{tempo-agendamento-e-aleatoriedade}

O simulador depende de um relógio global e de diferentes eventos
programados. A análise anterior identificou atualizações de horário a
cada dez ciclos e outra regra associada ao produto 240 ×
velocidade\_tempo; a conversão desses elementos para horas e dias
simulados precisa ser estabelecida a partir do código inteiro e
verificada em execuções controladas. A ordem de movimento, picada,
mortalidade, transmissão e atualização de estados também deverá ser
explicitada, pois pode alterar os resultados. (HKIOSHI, 2026)

A reprodutibilidade deve registrar versão do código, versão do GAMA,
arquivos GIS, parâmetros, configuração experimental e sementes. O
repositório apresenta um parâmetro configurável de semente e também uma
atribuição fixa de 12345.0 na inicialização. Antes de declarar que o
parâmetro da interface controla a execução, é necessário verificar qual
atribuição prevalece e como o gerador pseudoaleatório é inicializado.
(HKIOSHI, 2026)

\subsection{3.10 Inicialização, escalonamento populacional e dados de
entrada}\label{inicializauxe7uxe3o-escalonamento-populacional-e-dados-de-entrada}

A configuração examinada anteriormente define quantidade\_pessoas igual
a 1.191 e distribui agentes de acordo com regiões da cidade. O valor é
quantidade de agentes sintéticos, não a população real de São Paulo. Uma
simulação reduzida deve justificar pesos, densidades, formação de
contatos e interpretação das contagens. Preservar apenas percentuais
demográficos não demonstra equivalência epidemiológica entre a população
sintética e a real. (HKIOSHI, 2026)

A matriz de entradas deve relacionar cada arquivo a sua finalidade: base
territorial, atributos populacionais, parâmetros ambientais, valores de
vacinação e configuração de cenários. Sempre que uma entrada representar
hipótese experimental e não informação observada, isso será declarado
explicitamente. Os arquivos externos precisam ser acompanhados de
origem, licença, data e transformações, quando disponíveis. (Macal;
North, 2010; Grimm et al., 2020)

\subsection{3.11 Saídas do modelo e
indicadores}\label{sauxeddas-do-modelo-e-indicadores}

As saídas incluem contagens por estado humano e vetorial, população
total, mortes, tempo simulado e resultados persistidos em CSV, conforme
as estruturas já encontradas no repositório. Antes de utilizá-las em
gráficos epidemiológicos, será necessário verificar as identidades de
contagem, a interpretação de cada coluna, a escala temporal e a
consistência das unidades. Contagens de infecções simuladas não
equivalem automaticamente a casos notificados nos sistemas de
vigilância. (HKIOSHI, 2026)

Quadro 7 --- Verificações mínimas das saídas

\begin{longtable}[]{@{}
  >{\raggedright\arraybackslash}p{(\linewidth - 2\tabcolsep) * \real{0.5000}}
  >{\raggedright\arraybackslash}p{(\linewidth - 2\tabcolsep) * \real{0.5000}}@{}}
\toprule\noalign{}
\begin{minipage}[b]{\linewidth}\raggedright
\textbf{Indicador}
\end{minipage} & \begin{minipage}[b]{\linewidth}\raggedright
\textbf{Checagem necessária}
\end{minipage} \\
\midrule\noalign{}
\endhead
\bottomrule\noalign{}
\endlastfoot
Estados humanos & S + E + I + R ≤ população humana total, com explicação
para outras condições \\
Estados vetoriais & S + E + I = total vivo se os estados forem
exaustivos \\
Mortes & Não negatividade, monotonicidade se o indicador for
acumulado \\
Tempo & Intervalo entre linhas do CSV e correspondência com ciclos \\
Reprodutibilidade & Comparar resultados com parâmetros e sementes
idênticos \\
\end{longtable}

Fonte: elaboração própria.

\subsection{3.12 Hipóteses, limitações e
rastreabilidade}\label{hipuxf3teses-limitauxe7uxf5es-e-rastreabilidade}

As principais simplificações identificadas ou previstas incluem uso de
população sintética, representação limitada dos sorotipos, agregação
espacial, probabilidades parametrizadas e critérios simplificados de
deslocamento e condições ambientais. Cada hipótese deverá ser rastreável
aos arquivos GAML, à literatura que a sustenta ou à declaração de que é
uma opção experimental. Uma mesma variável pode ser biologicamente
significativa e, ao mesmo tempo, insuficientemente fundamentada para uso
preditivo. (HKIOSHI, 2026)

A análise crítica inicial ainda aponta possíveis fragilidades na
inicialização de estados vetoriais, no calendário de eventos, na
compatibilidade das unidades espaciais e no significado epidemiológico
de algumas transições. Essas observações são hipóteses de defeito a
confirmar em revisão e testes, não constatação de falha em todas as
execuções. Os testes e resultados serão apresentados no Capítulo 5.
(Macal; North, 2010; Grimm et al., 2020)

\subsection{3.13 Síntese do capítulo}\label{suxedntese-do-capuxedtulo-1}

A escolha do GAMA com recursos GIS é sustentada por sua adequação à
representação integrada de entidades, espaço e experimentação e pela
correspondência com os módulos existentes no repositório. O modelo
computacional foi apresentado em termos de entidades, estados,
interações, entradas e saídas, mantendo separadas as regras desejadas e
as evidências da implementação. O Capítulo 4 apresentará os componentes
GAML com trechos de código rastreáveis, detalhará o agendamento e
examinará a consistência das operações descritas. (HKIOSHI, 2026)

NOTA DE ATUALIZAÇÃO --- 8/11/2026: esta análise preserva evidências
históricas referentes à versão inspecionada em 30/10/2026.
Posteriormente, em 7/11/2026, os commits 94754bf e 621703c foram
associados à issue \#12, com relatos de mudanças em sementes, população
vetorial, inicialização e gravação dos resultados. O encerramento da
issue não substitui a confirmação por testes nem permite atribuir os
CSVs históricos ao novo código. As descrições de defeitos anteriores
devem ser lidas como achados da versão histórica, e não como prova de
que persistem na versão atual. (HKIOSHI, 2026)

\subsection{4.0 Atualização de implementação e rastreabilidade da
revisão de novembro de
2026}\label{atualizauxe7uxe3o-de-implementauxe7uxe3o-e-rastreabilidade-da-revisuxe3o-de-novembro-de-2026}

Em 7 de novembro de 2026 foram registrados no repositório os commits
94754bf e 621703c, vinculados à issue \#12. Segundo o comentário
publicado na issue, houve modificações no registro da semente e nos
nomes dos arquivos de saída, no uso do parâmetro quantidade\_mosquito,
na restrição dos estados iniciais dos mosquitos e na condição de
transmissão humano--vetor. Essas alterações devem ser descritas como
modificações declaradas no histórico até que o comportamento seja
confirmado por execução. Os requisitos de probabilidade de transmissão,
eficácia vacinal, reprodutibilidade e exportação por cenário e réplica
permanecem sujeitos a auditoria por testes objetivos. GitHub:
https://github.com/hkioshi/TccAbm/issues/12. (HKIOSHI, 2026)

\section{4 ARQUITETURA E IMPLEMENTAÇÃO DO
SIMULADOR}\label{arquitetura-e-implementauxe7uxe3o-do-simulador}

Este capítulo analisa a implementação existente no repositório público
TccAbm, desenvolvida em GAML para a plataforma GAMA. A descrição foi
baseada na leitura dos arquivos-fonte da ramificação main em 26 de
outubro de 2026. Por se tratar de análise estática do código, não se
atribui a ela o mesmo estatuto de teste executado, validação
epidemiológica ou medição de desempenho. A área geográfica declarada
para o experimento é o município de São Paulo; sua representação
computacional utiliza feições geográficas e uma agregação operacional em
cinco regiões. (HKIOSHI, 2026)

\subsection{4.1 Organização arquitetural e responsabilidades dos
módulos}\label{organizauxe7uxe3o-arquitetural-e-responsabilidades-dos-muxf3dulos}

O arquivo Main.gaml funciona como ponto de entrada do experimento
gráfico, importando explicitamente Mosquito.gaml, Pessoa.gaml,
Ambiente.gaml, Poca.gaml e globals.gaml. O componente global, por sua
vez, centraliza parâmetros, inicialização, agendamento, ambiente
geográfico e exportação de dados. Existe ainda um componente
Imaturo.gaml, importado por Poca.gaml, destinado às fases anteriores à
vida adulta. Essa separação por arquivos favorece a leitura das
responsabilidades, embora importações cruzadas exijam cuidado com o
acoplamento (TCCABM, 2026). (HKIOSHI, 2026)

A representação visual da arquitetura deverá ser inserida no arquivo
final somente após conferência da figura e de sua legenda. Os módulos
efetivamente descritos no texto são Main.gaml, globals.gaml e as
espécies de agentes. (HKIOSHI, 2026)

\begin{longtable}[]{@{}
  >{\raggedright\arraybackslash}p{(\linewidth - 2\tabcolsep) * \real{0.5000}}
  >{\raggedright\arraybackslash}p{(\linewidth - 2\tabcolsep) * \real{0.5000}}@{}}
\toprule\noalign{}
\begin{minipage}[b]{\linewidth}\raggedright
\textbf{Arquivo}
\end{minipage} & \begin{minipage}[b]{\linewidth}\raggedright
\textbf{Responsabilidade identificada}
\end{minipage} \\
\midrule\noalign{}
\endhead
\bottomrule\noalign{}
\endlastfoot
Main.gaml & Parâmetros da interface, exibição cartográfica e
monitores \\
globals/globals.gaml & Inicialização, clima, tempo, população,
mortalidade e relatório CSV \\
agents/Pessoa.gaml & Atributos humanos, estados epidemiológicos,
mobilidade e natalidade \\
agents/Mosquito.gaml & Atributos vetoriais, picada, transmissão,
deslocamento e mortalidade \\
agents/Poca.gaml & Capacidade de criadouro conforme chuva e
temperatura \\
agents/Imaturo.gaml & Estágios ovo, larva e pupa e criação de adultos \\
agents/Ambiente.gaml & Feições territoriais, atributo de região e
representação gráfica \\
\end{longtable}

\subsection{4.2 Importação GIS e delimitação
municipal}\label{importauxe7uxe3o-gis-e-delimitauxe7uxe3o-municipal}

No componente global, o arquivo ../Gis/subprefeitura\_v2.shp é declarado
como entrada e seu envelope é utilizado para definir a geometria global
do mundo. A criação dos agentes territoriais emprega o atributo
nm\_regiao\_. A abordagem permite associar posições e agentes humanos a
geometrias reais, mas o nome do arquivo ou da variável não comprova por
si só a exatidão da fronteira municipal, do sistema de referência ou da
integridade topológica. Essas propriedades exigem inspeção do Shapefile
e de sua documentação cartográfica. (HKIOSHI, 2026)

\begin{longtable}[]{@{}
  >{\raggedright\arraybackslash}p{(\linewidth - 0\tabcolsep) * \real{1.0000}}@{}}
\toprule\noalign{}
\begin{minipage}[b]{\linewidth}\raggedright
\textbf{file subprefeituras\_shapefile \textless-
file("../Gis/subprefeitura\_v2.shp");\\
geometry shape \textless- envelope(subprefeituras\_shapefile);\\
create ambiente from: subprefeituras\_shapefile\\
with: {[}nm\_regiao\_::read("nm\_regiao\_"){]};}\strut
\end{minipage} \\
\midrule\noalign{}
\endhead
\bottomrule\noalign{}
\endlastfoot
\end{longtable}

Código --- fragmentos de globals.gaml, linhas 58--59 e 113 (TccAbm,
2026).

A implementação analisada utiliza cinco rótulos de região: Centro,
Norte, Leste, Sul e Oeste. Esse agrupamento é uma simplificação
operacional e não deve ser confundido com os distritos administrativos
ou com a totalidade das subprefeituras do município. O posicionamento
utiliza any\_location\_in sobre geometrias das regiões, devendo ser
conferidos o sistema de coordenadas, a unidade efetiva das distâncias e
a forma de tratar fronteiras. (HKIOSHI, 2026)

\subsection{4.3 Inicialização das populações
sintéticas}\label{inicializauxe7uxe3o-das-populauxe7uxf5es-sintuxe9ticas}

O parâmetro quantidade\_pessoas é fixado em 1.191, acompanhado de
comentário que o descreve como escala 1:10.000. A implementação também
cria 100 mosquitos na inicialização. Os humanos são alocados entre as
cinco regiões mediante fatores proporcionais expressos diretamente no
código; o total da região Oeste é calculado por diferença para preservar
o total de humanos definido. Os fatores efetivamente programados devem
prevalecer sobre percentuais descritos em comentários antigos, pois não
são idênticos. (HKIOSHI, 2026)

\begin{longtable}[]{@{}
  >{\raggedright\arraybackslash}p{(\linewidth - 0\tabcolsep) * \real{1.0000}}@{}}
\toprule\noalign{}
\begin{minipage}[b]{\linewidth}\raggedright
\textbf{int quantidade\_pessoas \textless- 1191;\\
int qtd\_centro \textless- int(quantidade\_pessoas * 0.0370);\\
int qtd\_norte \textless- int(quantidade\_pessoas * 0.1929);\\
int qtd\_leste \textless- int(quantidade\_pessoas * 0.3485);\\
int qtd\_sul \textless- int(quantidade\_pessoas * 0.3262);\\
int qtd\_oeste \textless- quantidade\_pessoas - qtd\_centro - qtd\_norte
- qtd\_leste - qtd\_sul;}\strut
\end{minipage} \\
\midrule\noalign{}
\endhead
\bottomrule\noalign{}
\endlastfoot
\end{longtable}

Código --- globals.gaml, linhas 62 e 128--132 (TccAbm, 2026).

Atribuir proporções regionais a uma amostra populacional não assegura,
isoladamente, equivalência de contatos e densidade com a população real
de São Paulo. A calibração espacial e a relação entre as grandezas
simuladas e observadas precisam ser tratadas na seção experimental.
(HKIOSHI, 2026)

\subsection{4.4 Agentes humanos e mobilidade
espacial}\label{agentes-humanos-e-mobilidade-espacial}

A espécie pessoa usa a skill moving e possui idade, condição vacinal,
estado epidemiológico, variáveis de duração da infecção, destino,
velocidade e atributos de mortalidade. A ação andar implementa
deslocamentos em janelas horárias entre 5h e 8h e entre 16h e 19h;
nessas janelas o destino é sorteado conforme uma tabela de
probabilidades por região e executado com goto. Fora delas, o indivíduo
realiza um movimento local por wander. A rotina também reposiciona o
agente quando ele é detectado fora das geometrias territoriais (TCCABM,
2026). (HKIOSHI, 2026)

\begin{longtable}[]{@{}
  >{\raggedright\arraybackslash}p{(\linewidth - 0\tabcolsep) * \real{1.0000}}@{}}
\toprule\noalign{}
\begin{minipage}[b]{\linewidth}\raggedright
\textbf{if ((hora \textgreater= 5 and hora \textless{} 8) or (hora
\textgreater= 16 and hora \textless{} 19)) \{\\
if (destino = nil) \{\\
destino \textless- sortear\_destino(location);\\
velocidade \textless- velocidade\_andando;\\
\}\\
if (destino != nil) \{\\
do goto target: destino speed: velocidade;\\
\}\\
\} else \{\\
do wander speed: velocidade amplitude: 360.0;\\
\}}\strut
\end{minipage} \\
\midrule\noalign{}
\endhead
\bottomrule\noalign{}
\endlastfoot
\end{longtable}

Código --- excerto abreviado de Pessoa.gaml, linhas 62--87 (TccAbm,
2026).

A rotina de escolha de destinos define distribuições distintas para
manhã e tarde. Não se identificou, nesses trechos, cálculo de caminhos
sobre um grafo viário; o movimento documentado deve ser descrito como
deslocamento para coordenadas espaciais, e não como roteamento
comprovado pela malha de ruas. É necessário verificar ainda se os
intervalos de seleção de destinos são completos e mutuamente exclusivos
para os valores sorteados. (HKIOSHI, 2026)

\subsection{4.5 Estados epidemiológicos
humanos}\label{estados-epidemioluxf3gicos-humanos}

O arquivo Pessoa.gaml estabelece a codificação 0 = suscetível, 1 =
exposto, 2 = infectado e 3 = recuperado. A ação MudarEstado incrementa o
estado e retorna ao suscetível após o estado 3. A progressão temporal
está na ação Envelhecer, que utiliza periodo\_incubacao, periodo\_viral
e periodo\_recuperacao. A condição vacinal é mantida separadamente, como
atributo booleano. Essa estrutura aproxima-se da especificação SEIR, mas
seu comportamento exato depende das condições de acionamento. (HKIOSHI,
2026)

Há um ponto de atenção: o contador dias\_Infectado só é incrementado
quando o estado é 1 ou 2. Como a condição de saída do estado 3 depende
desse contador atingir a soma dos três períodos, pode ocorrer
permanência indefinida em R caso ele não atinja o limite antes da
recuperação. O comportamento deverá ser demonstrado por teste unitário
de transições, não presumido a partir do nome do parâmetro. (HKIOSHI,
2026)

\subsection{4.6 Agente mosquito, transmissão e
reprodução}\label{agente-mosquito-transmissuxe3o-e-reproduuxe7uxe3o}

A espécie mosquito define sexo, idade, velocidade, posição associada à
poça natal, distância máxima, alcance de contato e parâmetros
probabilísticos. A ação Picar é convocada pelo agendamento global
somente para mosquitos identificados como fêmeas. Ela efetua o sorteio
de tentativa de picada, seleciona uma pessoa dentro da distância de
1.800 unidades espaciais e avalia combinações de estados que podem
provocar transmissão. A eficácia vacinal aparece como redução
multiplicativa da probabilidade efetiva de infecção do humano. (HKIOSHI,
2026)

\begin{longtable}[]{@{}
  >{\raggedright\arraybackslash}p{(\linewidth - 0\tabcolsep) * \real{1.0000}}@{}}
\toprule\noalign{}
\begin{minipage}[b]{\linewidth}\raggedright
\textbf{if (alvo.estado = 0 and estado = 2) \{\\
ask alvo \{\\
float prob\_efetiva \textless- prob\_transmissao;\\
if (vacinado) \{\\
prob\_efetiva \textless- prob\_transmissao * (1 - eficacia / 100.0);\\
\}\\
if (rnd(1, 100) \textless= prob\_efetiva) \{ do MudarEstado; \}\\
\}\\
\}}\strut
\end{minipage} \\
\midrule\noalign{}
\endhead
\bottomrule\noalign{}
\endlastfoot
\end{longtable}

Código --- excerto de Mosquito.gaml, linhas 81--92 (TccAbm, 2026).

O uso de rnd(0,1) na tentativa de picada e de rnd(1,100) na transmissão
exige explicitar a escala de cada parâmetro. O código declara
prob\_picada = 0,25 e prob\_transmissao = 27,50 no mosquito, enquanto o
agente humano apresenta outro parâmetro de transmissão. Não é
metodologicamente correto tratá-los como a mesma probabilidade. Também
há condição que considera humanos expostos ou infectados aptos a
transmitir ao mosquito, aspecto que deve ser justificado biologicamente
ou ajustado. (HKIOSHI, 2026)

A rotina de reprodução cria um lote de 100 agentes imaturos ao encontrar
poça próxima e capacidade residual. Como a condição antecedente apenas
testa se a quantidade atual é menor que a capacidade, deve-se verificar
se a criação do lote pode ultrapassar o limite. O vínculo entre
reprodução e a existência de pessoa próxima é outra hipótese da
implementação que demanda revisão. (HKIOSHI, 2026)

\subsection{4.7 Desenvolvimento imaturo e dinâmica de
criadouros}\label{desenvolvimento-imaturo-e-dinuxe2mica-de-criadouros}

Imaturo.gaml implementa três estágios numéricos: ovo, larva e pupa. A
duração de cada fase é sorteada em função da temperatura global; quando
o estágio 2 termina, o agente imaturo cria um mosquito adulto na
localização da poça e é removido. Poca.gaml representa a disponibilidade
hídrica por um atributo tamanho, limitado a cinco níveis. A intensidade
de chuva aumenta o nível e a ausência de chuva o diminui, enquanto a
capacidade é expressa pelo produto tamanho × 20. Essas relações são
hipóteses operacionais e não equivalem, automaticamente, a uma medição
entomológica no município. (Taillandier et al., 2019; GAMA Platform,
{[}s. d.{]})

A leitura estática identifica o mecanismo de atualização ambiental, mas
não comprova que todas as listas de imaturos estejam sempre
sincronizadas após mortalidade e transformação em adultos. Essa
invariável deve ser testada. Também é importante documentar se as
durações sorteadas em dias são aplicadas na mesma unidade de Envelhecer,
o que depende do agendamento global. (HKIOSHI, 2026)

\subsection{4.8 Controle temporal, clima e variabilidade
estocástica}\label{controle-temporal-clima-e-variabilidade-estocuxe1stica}

Na versão atual de globals.gaml, a atualização da hora está configurada
para ocorrer a cada cinco ciclos (tempo\_hora = 5), e o evento diário a
cada 120 × velocidade\_tempo ciclos (tempo\_dia = 24 × 5). Esses
parâmetros diferem da versão anteriormente examinada. O modelo também
emprega atualização climática e eventos de agentes, cuja ordem de
execução deve ser verificada no GAMA. (Taillandier et al., 2019; GAMA
Platform, {[}s. d.{]})

Registro da auditoria histórica (versão de outubro de 2026): A interface
expõe o parâmetro seed, mas a rotina init atribui seed \textless-
12345.0. Trata-se de um risco de reprodutibilidade a ser verificado em
execução, principalmente nos experimentos batch que declaram 30
repetições. Observação: a issue \#12 relata alterações posteriores em
7/11/2026; confirmar a situação atual no código e por testes antes de
repetir esta crítica como defeito vigente. (HKIOSHI, 2026)

\subsection{4.9 Interface, indicadores e persistência das
saídas}\label{interface-indicadores-e-persistuxeancia-das-sauxeddas}

O experimento MeuExperimento, em Main.gaml, disponibiliza velocidade do
tempo, percentual de vacinados, semente e eficácia vacinal. A
visualização bidimensional exibe ambiente, pessoas, mosquitos e poças.
Monitores mostram ciclo, dia, mês, hora, população por estado e mortes
por dengue. O arquivo globals.gaml contém a ação Relatorio, que grava em
resultados.csv totais por classe, estados epidemiológicos e mortalidade
por faixas etárias. (HKIOSHI, 2026)

A estrutura de gravação é adequada para produzir uma série temporal de
indicadores. Contudo, a ausência de campos explícitos para a versão do
código, conjunto de parâmetros, semente efetiva e identificador de
repetição dificulta a auditoria de experimentos. Recomenda-se um
manifesto por execução e verificações automáticas de consistência:
totais de humanos e mosquitos devem corresponder à soma das respectivas
categorias válidas, desde que as categorias sejam mutuamente exclusivas
e exaustivas. (HKIOSHI, 2026)

\subsection{4.10 Revisão crítica: divergências e requisitos de
verificação}\label{revisuxe3o-cruxedtica-diverguxeancias-e-requisitos-de-verificauxe7uxe3o}

\begin{longtable}[]{@{}
  >{\raggedright\arraybackslash}p{(\linewidth - 4\tabcolsep) * \real{0.3333}}
  >{\raggedright\arraybackslash}p{(\linewidth - 4\tabcolsep) * \real{0.3333}}
  >{\raggedright\arraybackslash}p{(\linewidth - 4\tabcolsep) * \real{0.3333}}@{}}
\toprule\noalign{}
\begin{minipage}[b]{\linewidth}\raggedright
\textbf{Achado na leitura estática}
\end{minipage} & \begin{minipage}[b]{\linewidth}\raggedright
\textbf{Risco ou interpretação}
\end{minipage} & \begin{minipage}[b]{\linewidth}\raggedright
\textbf{Teste ou ação recomendada}
\end{minipage} \\
\midrule\noalign{}
\endhead
\bottomrule\noalign{}
\endlastfoot
rnd(0,3) na inicialização dos mosquitos & Possível estado 3 não
monitorado, se os limites forem inclusivos & Assegurar domínio \{0,1,2\}
e testar toda inicialização \\
MudarEstado do mosquito com mod 3 & Possível retorno I → S em
acionamentos adicionais & Verificar ciclo de vida epidemiológico e
trajetórias \\
Contador humano não cresce no estado R & Possível impossibilidade de R →
S & Teste das transições com relógio controlado \\
seed definida na interface e no init & Semente da interface pode ser
sobrescrita & Testar igualdade e variação entre execuções \\
240 × velocidade\_tempo para evento diário & Inconsistência potencial de
escala temporal & Tabela de equivalência ciclo--hora--dia \\
Lote de 100 imaturos com teste apenas do limite inicial & Possível
excesso de capacidade da poça & Teste com capacidade residual de 1 a
99 \\
Distâncias 1.800 e 50.000 sem CRS documentado & Impossibilidade de
interpretar alcances fisicamente & Documentar projeção e unidade
espacial \\
População sintética de 1.191 pessoas & Resultados não equivalentes
automaticamente à cidade real & Documentar escala, densidade e
sensibilidade \\
\end{longtable}

Esses achados não são equivalentes a erros confirmados em execução;
descrevem riscos, inconsistências textuais ou consequências possíveis da
lógica examinada. A etapa seguinte deverá gerar evidências por meio de
testes reproduzíveis e registros de execução, com identificação de
versão e parâmetros. (HKIOSHI, 2026)

\subsection{4.11 Configuração de experimentos em lote
(batch)}\label{configurauxe7uxe3o-de-experimentos-em-lote-batch}

A versão pública de Main.gaml consultada em 30 de outubro de 2026
contém, além do experimento interativo MeuExperimento, quatro definições
de experimentos do tipo batch. Na linguagem GAML, essas declarações
especificam variações de parâmetros, números de repetições e condições
de término. Sua existência no código comprova a configuração
experimental, mas não a conclusão das execuções nem a integridade das
saídas. (HKIOSHI, 2026)

Quadro 4.1 --- Experimentos declarados em Main.gaml. Fonte: elaboração
própria a partir de TccAbm (2026), Main.gaml, linhas 71--129.

E03Variabilidade --- repeat: 30; porcentagem\_de\_vacinados = 0;
eficacia = 0; condição de término: cycle \textgreater= 87600. (HKIOSHI,
2026)

E04Vacinacao --- repeat: 30; porcentagem\_de\_vacinados de 0 a 100,
passo 25; eficacia = 100; condição de término: cycle \textgreater=
87600. (HKIOSHI, 2026)

E05PopulacaoVetorial --- repeat: 30; quantidade\_mosquito de 50 a 400,
passo 50; vacinação desativada; condição de término: cycle \textgreater=
87600. (HKIOSHI, 2026)

E07Escalabilidade --- repeat: 1; quantidade\_pessoas de 500 a 10.000,
passo 500; quantidade\_mosquito = 100; vacinação desativada; condição de
término: cycle \textgreater= 87600. (HKIOSHI, 2026)

Uma distinção arquitetural necessária é que o número de repetições
indicado por repeat não substitui o registro de identificadores de
execução e sementes efetivamente utilizadas. O experimento de
escalabilidade tem apenas uma execução configurada por combinação; por
isso, sua definição atual não basta para caracterizar estatisticamente a
variabilidade do desempenho. (HKIOSHI, 2026)

\subsection{4.12 Contrato entre parâmetros, inicialização e
execução}\label{contrato-entre-paruxe2metros-inicializauxe7uxe3o-e-execuuxe7uxe3o}

Registro da auditoria histórica (versão de outubro de 2026): A
configuração externa de um parâmetro só se torna efetiva se as rotinas
de inicialização e atualização o consumirem. Em globals.gaml, embora
quantidade\_mosquito seja declarada com valor inicial 100 e variada no
experimento E05PopulacaoVetorial, o bloco de inicialização continua
utilizando create mosquito number: 100. Portanto, na leitura estática do
código, não se identifica aplicação direta da variável
quantidade\_mosquito na criação inicial. A alteração da variável no
batch não permite, por si só, afirmar que o número inicial de mosquitos
mudou. Observação: a issue \#12 relata alterações posteriores em
7/11/2026; confirmar a situação atual no código e por testes antes de
repetir esta crítica como defeito vigente. (HKIOSHI, 2026)

Registro da auditoria histórica (versão de outubro de 2026): O parâmetro
seed é exposto no experimento gráfico, mas a rotina global init executa
seed \textless- 12345.0. Esta atribuição merece teste controlado,
especialmente nas 30 replicações, para determinar como o mecanismo
interno de sementes do GAMA e a atribuição no código interagem. Sem
registrar a semente efetiva de cada execução, não é possível assegurar
independência estatística nem reprodução exata dos resultados.
Observação: a issue \#12 relata alterações posteriores em 7/11/2026;
confirmar a situação atual no código e por testes antes de repetir esta
crítica como defeito vigente. (HKIOSHI, 2026)

O mecanismo temporal usa tempo\_hora = 5 ciclos e tempo\_dia = 24 × 5 =
120 ciclos. A atualização de hora ocorre a cada 5 ciclos, enquanto o
evento diário usa 120 × velocidade\_tempo ciclos. Assim, a
compatibilidade entre hora, dia, velocidade\_tempo e condição de término
cycle \textgreater= 87600 precisa ser verificada para cada cenário. Com
velocidade\_tempo = 1, 87.600/120 = 730 intervalos diários de 120
ciclos; a duração real da série registrada depende também da lógica de
atualização do contador dia e dos demais eventos. (HKIOSHI, 2026)

\subsection{4.13 Registro de resultados, rastreabilidade e requisitos de
verificação}\label{registro-de-resultados-rastreabilidade-e-requisitos-de-verificauxe7uxe3o}

A ação Relatorio grava 19 campos em resultados.csv: dia, chuva,
temperatura, quatro indicadores humanos mais o total, três
compartimentos vetoriais mais o total, e sete contadores de mortalidade.
O cabeçalho é regravado na inicialização (rewrite: true), enquanto as
linhas de resultado são acrescentadas posteriormente (rewrite: false).
Este desenho fornece uma série temporal para uma execução, mas pode
apagar registros anteriores e não inclui automaticamente o nome do
experimento, a repetição, a semente, os parâmetros, o commit ou o tempo
de processamento. (HKIOSHI, 2026)

Registro da auditoria histórica (versão de outubro de 2026): A
representação dos mosquitos utiliza contagens para estados 0, 1 e 2;
entretanto, a inicialização executa estado \textless- rnd(0,3), que
também admite o valor 3. Além disso, a ação MudarEstado usa (estado + 1)
mod 3. Estas construções exigem verificar os invariantes de contagem e
as transições esperadas para a progressão vetorial antes de interpretar
as séries epidemiológicas. Observação: a issue \#12 relata alterações
posteriores em 7/11/2026; confirmar a situação atual no código e por
testes antes de repetir esta crítica como defeito vigente. (HKIOSHI,
2026)

Requisitos de verificação resultantes da auditoria estática: (a)
assegurar que todas as entidades estejam em estados válidos e que a soma
dos compartimentos corresponda aos totais; (b) comprovar que os
parâmetros alterados pelo batch mudam o comportamento inicial ou
dinâmico esperado; (c) guardar semente e identificador por réplica; (d)
aferir a duração temporal real e a frequência de registro; (e) separar
resultados por execução; (f) medir explicitamente tempo e memória para
sustentar alegações de escalabilidade. A realização desses testes
pertence à avaliação experimental do Capítulo 5. (HKIOSHI, 2026)

Versão consultada: branch main do repositório público TccAbm, em 30 de
outubro de 2026. A inspeção foi estática, a partir dos arquivos
Main.gaml, globals.gaml e Mosquito.gaml disponibilizados publicamente;
não foi efetuada execução da plataforma GAMA nesta revisão. (HKIOSHI,
2026)

\subsection{4.14 Síntese do capítulo}\label{suxedntese-do-capuxedtulo-2}

A implementação examinada reúne módulos GAML, agentes humanos e
vetoriais, dados GIS, estados imaturos, atualização ambiental e
mecanismos de exportação. A versão mais recente acrescenta quatro
experimentos batch. Sua descrição técnica, porém, revela contratos de
parâmetros e invariantes epidemiológicos que demandam verificação
funcional antes de interpretar resultados. A avaliação experimental
correspondente será conduzida no Capítulo 5. (Taillandier et al., 2019;
GAMA Platform, {[}s. d.{]})

NOTA DE ATUALIZAÇÃO --- 8/11/2026: esta análise preserva evidências
históricas referentes à versão inspecionada em 30/10/2026.
Posteriormente, em 7/11/2026, os commits 94754bf e 621703c foram
associados à issue \#12, com relatos de mudanças em sementes, população
vetorial, inicialização e gravação dos resultados. O encerramento da
issue não substitui a confirmação por testes nem permite atribuir os
CSVs históricos ao novo código. As descrições de defeitos anteriores
devem ser lidas como achados da versão histórica, e não como prova de
que persistem na versão atual. (HKIOSHI, 2026)

\subsection{5.0 Alcance das evidências e versões
analisadas}\label{alcance-das-eviduxeancias-e-versuxf5es-analisadas}

As 229 linhas históricas atribuídas ao CSV consultado em outubro de 2026
constituem registros observacionais da versão anterior à revisão de
7/11/2026. Elas não são evidência de que E03, E04 ou E05 tenham sido
executados no código atualizado, nem comprovam adequação das correções.
Os resultados que dependem de mudanças na seed, nas probabilidades de
transmissão, na população de mosquitos e na vacinação somente poderão
ser acrescentados depois de testados os commits correspondentes. Até
então, os resultados históricos devem ser identificados como análise
exploratória e auditoria de integridade. (HKIOSHI, 2026)

\section{5 VERIFICAÇÃO, EXPERIMENTAÇÃO, RESULTADOS E
DISCUSSÃO}\label{verificauxe7uxe3o-experimentauxe7uxe3o-resultados-e-discussuxe3o}

Este capítulo examina as evidências experimentais atualmente disponíveis
para o simulador baseado em agentes implementado na plataforma GAMA,
cuja área geográfica de aplicação é o município de São Paulo. São
examinadas saídas arquivadas no repositório TccAbm, definidos critérios
de verificação da implementação e estabelecido um protocolo para a
avaliação estocástica, epidemiológica e computacional. A análise
diferencia expressamente os resultados observáveis no arquivo de saída,
os procedimentos ainda não executados nesta revisão e as limitações que
restringem inferências sobre o comportamento real da dengue. (HKIOSHI,
2026)

Escopo da revisão de 30 out. 2026. As afirmações sobre o estado recente
do repositório referem-se ao commit f3b3a424; não houve execução local
do simulador GAMA. O CSV examinado não informa a identidade das
repetições e as figuras herdadas da revisão anterior mostram apenas um
trecho histórico. Não se declaram testes aprovados ou cenários
comparativos concluídos. (HKIOSHI, 2026)

A simples execução de um modelo não constitui demonstração de validade
epidemiológica. Verificação significa examinar a correspondência entre
código e especificação; validação investiga a adequação da representação
ao fenômeno e aos propósitos definidos; calibração ajusta parâmetros com
dados de referência; e avaliação de desempenho caracteriza o consumo de
recursos computacionais. Essas atividades precisam ser documentadas
separadamente (GRIMM et al., 2020; MACAL; NORTH, 2010).

\subsection{5.1 Fontes de evidência e
rastreabilidade}\label{fontes-de-eviduxeancia-e-rastreabilidade}

A fonte primária desta atualização é o arquivo resultados.csv no
repositório público TccAbm, examinado na revisão de 30 out. 2026 e
referenciado ao commit f3b3a42414129268888e2120a275b1bcd290d525. Essa
versão contém 229 registros de dados e 19 colunas, incluindo dia, chuva,
temperatura, contagens humanas, estados dos mosquitos e mortalidade. O
nome do arquivo e a estrutura das colunas não identificam, por si
mesmos, qual experimento produziu cada linha. A branch principal do
repositório pode receber modificações posteriores: por isso, as
constatações abaixo se referem ao instantâneo explicitamente
identificado. (HKIOSHI, 2026)

O arquivo CSV, isoladamente, não informa o identificador exato do commit
executado, a versão do GAMA, a semente efetivamente utilizada, as
dependências ambientais, as opções do experimento, o hardware e as
condições de coleta. Portanto, os resultados são tratados neste capítulo
como saídas arquivadas e exploratórias, sem afirmar que tenham sido
reproduzidos nesta análise. (HKIOSHI, 2026)

Quadro 5.1 --- Fontes utilizadas e alcance de suas evidências

\begin{longtable}[]{@{}
  >{\raggedright\arraybackslash}p{(\linewidth - 4\tabcolsep) * \real{0.3333}}
  >{\raggedright\arraybackslash}p{(\linewidth - 4\tabcolsep) * \real{0.3333}}
  >{\raggedright\arraybackslash}p{(\linewidth - 4\tabcolsep) * \real{0.3333}}@{}}
\toprule\noalign{}
\begin{minipage}[b]{\linewidth}\raggedright
\textbf{Fonte}
\end{minipage} & \begin{minipage}[b]{\linewidth}\raggedright
\textbf{Permite verificar}
\end{minipage} & \begin{minipage}[b]{\linewidth}\raggedright
\textbf{Não demonstra isoladamente}
\end{minipage} \\
\midrule\noalign{}
\endhead
\bottomrule\noalign{}
\endlastfoot
Código GAML no GitHub & Regras e organização da implementação & Execução
correta ou validade biológica \\
resultados.csv & Séries numéricas e verificações aritméticas &
Identificação da execução, replicações ou calibração \\
Garcias (2024); Oliveira et al. (2025) & Abordagens e experimentos de
referência & Resultados obtidos pelo simulador de São Paulo \\
Bases públicas de vigilância (a selecionar) & Observações da dengue em
São Paulo & Equivalência direta com infecções simuladas \\
\end{longtable}

Fonte: elaboração própria (2026).

\subsection{5.2 Configuração da execução e limites das informações
disponíveis}\label{configurauxe7uxe3o-da-execuuxe7uxe3o-e-limites-das-informauxe7uxf5es-disponuxedveis}

No primeiro registro do CSV analisado, observam-se 1.191 agentes humanos
e 100 mosquitos. As contagens humanas são 1.189 suscetíveis, 2 expostos,
0 infectados e 0 recuperados, o que totaliza 1.191. Já as contagens de
mosquitos são 28 suscetíveis, 29 expostos e 21 infectados, totalizando
78, e não os 100 indicados na coluna mosquitos\_total. Trata-se de
divergência verificável entre indicadores exportados, que exige
auditoria do estado inicial e da lógica de contabilização. (HKIOSHI,
2026)

Embora o CSV apresente 229 linhas de dados, a coluna dia não constitui
um identificador inequívoco de execução: após o registro de dia 30, a
numeração retorna para 1. Sem identificador de execução, semente e
cenário, não é possível afirmar quantas execuções independentes estão
contidas no arquivo, nem usar suas linhas como uma única série
cronológica contínua. Antes da análise estatística, o CSV deve ser
segmentado e vinculado a metadados de execução. (HKIOSHI, 2026)

Quadro 5.2 --- Metadados necessários para um experimento reprodutível

\begin{longtable}[]{@{}
  >{\raggedright\arraybackslash}p{(\linewidth - 2\tabcolsep) * \real{0.5000}}
  >{\raggedright\arraybackslash}p{(\linewidth - 2\tabcolsep) * \real{0.5000}}@{}}
\toprule\noalign{}
\begin{minipage}[b]{\linewidth}\raggedright
\textbf{Elemento}
\end{minipage} & \begin{minipage}[b]{\linewidth}\raggedright
\textbf{Situação nesta análise}
\end{minipage} \\
\midrule\noalign{}
\endhead
\bottomrule\noalign{}
\endlastfoot
Commit e versão integral dos dados & A registrar no protocolo
definitivo \\
Versão da plataforma GAMA & A confirmar na máquina da execução \\
Semente aleatória realmente utilizada & A confirmar por
instrumentação \\
Duração e unidade temporal & Há dias rotulados de 0 a 14; sem
confirmação externa do calendário \\
Hardware, sistema operacional e memória & Não informado no CSV \\
Quantidade de replicações & Não demonstrada pelo arquivo único \\
Resolução e projeção do mapa GIS & A documentar a partir das bases
efetivamente carregadas \\
\end{longtable}

Fonte: elaboração própria (2026).

\subsection{5.3 Resultados exploratórios presentes no
repositório}\label{resultados-exploratuxf3rios-presentes-no-reposituxf3rio}

\subsubsection{5.3.1 Evolução dos estados
humanos}\label{evoluuxe7uxe3o-dos-estados-humanos}

A Figura 5.1 foi gerada em versão anterior do estudo a partir de um
trecho inicial de 15 registros; por isso, deve ser lida como ilustração
exploratória de um subconjunto histórico, e não como representação dos
229 registros agora identificados. Na versão atual do CSV, entre os dias
rotulados 0 e 30 do primeiro segmento, os infectados humanos passam de 0
a 382, ao passo que os suscetíveis caem de 1.189 a 121. Esses valores
descrevem saídas registradas, não casos reais de dengue nem desempenho
epidemiológico validado. (HKIOSHI, 2026)

Figura 5.1 --- Curvas exploratórias de estados humanos em subconjunto
histórico de 15 linhas. Fonte: elaboração própria a partir de registro
anterior de resultados.csv; não representa todos os dados do instantâneo
de 30 out. 2026. {[}Figura histórica: conferir dados e vínculo ao commit
antes de manter na versão final.{]}

O teste de integridade dos estados humanos deve ser reaplicado
sistematicamente a todos os 229 registros: humanos\_total = suscetíveis
+ expostos + infectados + recuperados, com exame separado de mortes,
nascimentos e eventuais categorias adicionais. A concordância em um
registro não garante a consistência da série inteira. Registros que
divergirem devem ser rastreados até o instante de atualização, a lógica
de transição e as operações de escrita no CSV. (HKIOSHI, 2026)

\subsubsection{5.3.2 Evolução e consistência da população
vetorial}\label{evoluuxe7uxe3o-e-consistuxeancia-da-populauxe7uxe3o-vetorial}

A Figura 5.2 documenta divergências observadas na amostra histórica
usada na versão anterior. Na versão atual do CSV, a diferença já aparece
no registro inicial: mosquitos\_total = 100, mas mosquitos\_suscetiveis
+ mosquitos\_expostos + mosquitos\_infectados = 78, diferença de 22
indivíduos. A leitura do código sugere como hipótese a inicialização de
um estado adicional ao conjunto contabilizado; essa explicação permanece
condicionada a teste de execução e inspeção dos agentes. (HKIOSHI, 2026)

Figura 5.2 --- Divergência exploratória entre totais e estados vetoriais
em subconjunto histórico. Fonte: elaboração própria a partir de registro
anterior de resultados.csv; requer reavaliação integral. {[}Figura
histórica: conferir dados e vínculo ao commit antes de manter na versão
final.{]}

Uma hipótese plausível é a existência de indivíduos em um estado não
contabilizado ou de diferenças no momento de atualização dos monitores.
Entretanto, nenhuma dessas hipóteses pode ser apresentada como causa
comprovada sem inspecionar a inicialização dos estados, os filtros dos
contadores e o escalonamento das exportações. A divergência é um achado
de consistência a ser investigado. (HKIOSHI, 2026)

\subsubsection{5.3.3 Indicadores de
mortalidade}\label{indicadores-de-mortalidade}

No primeiro segmento disponível do CSV recente, entre as linhas
rotuladas 0 e 30, os campos de mortalidade por dengue aparecem com valor
zero. Mais adiante no arquivo existem registros de mortalidade não nula.
Sem uma segmentação confiável das execuções, não é adequado interpretar
esses valores como uma curva única de óbitos ou calcular uma taxa de
letalidade agregada. (HKIOSHI, 2026)

\subsection{5.4 Verificação funcional: critérios e casos de
teste}\label{verificauxe7uxe3o-funcional-crituxe9rios-e-casos-de-teste}

Seguindo a separação metodológica entre verificação e validação
discutida por Arifin, Madey e Collins (2016), a primeira etapa dos
testes deve examinar invariantes computacionais, como conservação e
exclusividade dos estados, consistência dos totais registrados, controle
das sementes e coerência entre parâmetros batch e número de agentes
inicializados. Somente após essas verificações caberá uma análise de
ajuste ou plausibilidade epidemiológica; a documentação científica de
outros simuladores não substitui evidências produzidas no modelo deste
trabalho.

A verificação deve conectar cada requisito da especificação aos testes
das ações implementadas. Propõe-se uma bateria de testes unitários e
cenários controlados, seguida por testes de integração. Para cada teste,
será necessário registrar entradas, semente, resultado esperado,
resultado observado, evidência (log, captura de tela ou CSV) e situação
de aprovação. (HKIOSHI, 2026)

Quadro 5.3 --- Matriz de verificação proposta

\begin{longtable}[]{@{}
  >{\raggedright\arraybackslash}p{(\linewidth - 4\tabcolsep) * \real{0.3333}}
  >{\raggedright\arraybackslash}p{(\linewidth - 4\tabcolsep) * \real{0.3333}}
  >{\raggedright\arraybackslash}p{(\linewidth - 4\tabcolsep) * \real{0.3333}}@{}}
\toprule\noalign{}
\begin{minipage}[b]{\linewidth}\raggedright
\textbf{Identificação}
\end{minipage} & \begin{minipage}[b]{\linewidth}\raggedright
\textbf{Teste e condição}
\end{minipage} & \begin{minipage}[b]{\linewidth}\raggedright
\textbf{Critério esperado}
\end{minipage} \\
\midrule\noalign{}
\endhead
\bottomrule\noalign{}
\endlastfoot
V01 & Inicialização do mosquito & Estado inicial pertencente ao domínio
epidemiológico definido \\
V02 & Progressão dos estados vetoriais & Sem retorno não previsto de
infectante para suscetível \\
V03 & Contagem dos estados & Soma das classes compatível com total e
eventuais estados adicionais explicitados \\
V04 & Transmissão sem contato & Nenhuma nova infecção quando contato
impossível \\
V05 & Transmissão com probabilidade zero & Nenhuma transmissão atribuída
ao evento \\
V06 & Vacinação & Efeito corresponde à regra implementada e ao
parâmetro \\
V07 & Relógio & Conversão consistente entre ciclo, hora e dia
simulados \\
V08 & Semente & Mesma configuração e semente produzem saídas
reproduzíveis sob ambiente controlado \\
V09 & GIS e mobilidade & Agentes permanecem em geometrias/rotas
admissíveis segundo o modelo \\
V10 & Exportação & Campos, unidades, timestamps e contadores
coerentes \\
\end{longtable}

Fonte: elaboração própria (2026). Testes propostos, ainda não executados
nesta revisão.

\subsection{5.5 Planejamento da análise estocástica e de
sensibilidade}\label{planejamento-da-anuxe1lise-estocuxe1stica-e-de-sensibilidade}

O repositório contém quatro definições de experimento em lote,
identificadas como E03Variabilidade, E04Vacinacao, E05PopulacaoVetorial
e E07Escalabilidade. Essa constatação altera o estado do trabalho: não
se trata apenas de planejamento abstrato, pois os cenários estão
declarados no código. Contudo, a presença de uma definição de
experimento não demonstra que suas execuções ocorreram, que os
parâmetros foram efetivamente utilizados ou que as saídas foram
arquivadas de maneira rastreável. O protocolo de validação deve
anteceder qualquer inferência comparativa. (HKIOSHI, 2026)

Uma estratégia de investigação consiste em variar isoladamente a
cobertura vacinal, a eficácia representada, a probabilidade de picada, a
mortalidade e a quantidade inicial de vetores. As saídas de interesse
incluem pico de infectados, tempo até o pico, infecções acumuladas,
estados vetoriais e tempo de execução. Deve-se distinguir a exploração
local, variando um parâmetro de cada vez, de uma análise global, que
considera interações entre fatores. Métodos como Morris e Sobol podem
ser discutidos, conforme viabilidade computacional (MORRIS, 1991; SOBOL,
2001). (Gamal et al., 2025; Morris, 1991; Sobol, 2001)

\subsubsection{5.5.1 Desenho experimental proposto: fatores, respostas e
replicações}\label{desenho-experimental-proposto-fatores-respostas-e-replicauxe7uxf5es}

Para justificar a seleção das métricas, Doege (2026) constitui um
exemplo de análise de estratégias de controle da dengue com múltiplos
cenários, repetição estocástica e acompanhamento de indicadores como
pico de infecção, tempo ao pico, total acumulado e duração da epidemia.
O desenho do presente estudo pode adotar indicadores comparáveis, desde
que suas definições sejam explicitadas e as respectivas séries temporais
sejam realmente geradas pelo GAMA. A presença dessas métricas na
literatura não comprova que os experimentos E03, E04 e E05 tenham sido
executados no repositório analisado.

A sistematização do planejamento experimental adota como inspiração
didática a comparação gradual de cenários apresentada por Sanderson
(2020), adaptando-a à transmissão vetorial da dengue. Os procedimentos a
seguir constituem um protocolo a executar, não resultados observados.
Antes da campanha, é necessário corrigir ou justificar as
inconsistências de contagem dos estados vetoriais, revisar a relação
entre ticks e tempo epidemiológico, e confirmar o funcionamento do
controle de sementes no código GAML.

O cenário de referência deverá ser descrito por uma configuração
versionada do modelo: commit do repositório, versão do GAMA, arquivo
geográfico e respectivo sistema de coordenadas, duração e passo
temporal, população sintética, número inicial de mosquitos, parâmetros
de transmissão, mobilidade, ambiente e critérios de encerramento. Cada
execução deverá receber um identificador único e preservar a
configuração efetivamente utilizada. (Oliveira et al., 2025; Garcias,
2024)

Recomenda-se iniciar por uma avaliação de variabilidade com uma
configuração fixa e várias sementes distintas, verificando a
estabilidade de médias, dispersões e intervalos de incerteza de
indicadores predefinidos. Para comparação entre cenários, sempre que
adequado, pode-se empregar o mesmo conjunto de sementes em cada condição
para reduzir ruído de comparação; deve-se ainda avaliar se essa
estratégia é compatível com a ordem de sorteios do simulador. O número
final de replicações dependerá da precisão obtida no estudo piloto e
será documentado. (Gamal et al., 2025; Morris, 1991; Sobol, 2001)

Em seguida, poderão ser contrastados, um fator por vez, níveis
previamente justificados de: (a) probabilidade de transmissão associada
à interação humano--vetor; (b) abundância inicial dos mosquitos; (c)
parâmetros da mobilidade humana; e (d) distribuição espacial inicial de
agentes e criadouros. A escolha das faixas deverá respeitar a
parametrização efetiva do modelo e as fontes epidemiológicas, evitando
percentuais arbitrários interpretados como realistas. Experimentos
multifatoriais poderão investigar interações após a etapa exploratória.
(HKIOSHI, 2026)

Como variáveis-resposta, sugere-se registrar infecções humanas
incidentes e acumuladas, pico de infectados, tempo até o pico, população
vetorial por estado e distribuição espacial dos eventos, desde que os
respectivos contadores sejam verificados. Para avaliação de engenharia
de software, registrar também duração da execução, consumo de memória e
número de agentes. Comparações devem apresentar valores centrais e
dispersão, bem como as condições que diferem entre cenários. (Gamal et
al., 2025; Morris, 1991; Sobol, 2001)

Um indicador de reprodução de infecções humanas só poderá ser estimado
de maneira epidemiologicamente defensável se a implementação registrar
eventos de transmissão e possibilitar sua atribuição ao ciclo
humano--vetor--humano, com definição explícita de geração e janela
temporal. Não se deve transplantar diretamente o cálculo de R ou R₀ da
demonstração SIR do vídeo para um modelo com transmissão vetorial.
(Oliveira et al., 2025; Garcias, 2024)

Por fim, observações sobre pontos de circulação compartilhados e viagens
entre comunidades podem orientar perguntas adicionais sobre mobilidade,
mas tais mecanismos devem estar efetivamente implementados e
documentados antes de serem transformados em experimentos. Nenhuma
conclusão causal sobre intervenções ou previsões para a cidade de São
Paulo decorre exclusivamente das animações ou dos exemplos do vídeo.
(HKIOSHI, 2026)

\subsubsection{5.5.2 Quantificação da incerteza e variabilidade entre
execuções}\label{quantificauxe7uxe3o-da-incerteza-e-variabilidade-entre-execuuxe7uxf5es}

A análise proposta por Gamal et al. (2025) distingue quatro fontes de
incerteza em modelos comportamentais: (i) incerteza paramétrica,
decorrente dos valores e das distribuições atribuídas às entradas; (ii)
incerteza estrutural, associada às simplificações e às hipóteses da
representação; (iii) variabilidade do conjunto de execuções (ensemble),
associada aos mecanismos estocásticos; e (iv) incerteza empírica,
relativa à observação e à qualidade dos dados do sistema real. Essas
fontes não devem ser confundidas com defeitos de implementação:
inconsistências lógicas e falhas de contagem precisam ser corrigidas ou
justificadas na verificação, antes da interpretação estatística das
saídas.

Como etapa inicial, propõe-se executar um experimento piloto com
parâmetros, dados geográficos, quantidade de agentes e horizonte
temporal constantes, alterando somente a semente pseudoaleatória. Para
cada execução r e instante t, devem ser registrados os indicadores de
interesse y(r,t), permitindo estimar média amostral, desvio-padrão,
quantis empíricos e intervalos de Monte Carlo. O número definitivo de
replicações não será fixado sem evidência: deverá ser ampliado até
atingir precisão previamente estabelecida para as medidas escolhidas,
observando autocorrelação temporal e eventuais trajetórias com extinção
precoce. (Gamal et al., 2025; Morris, 1991; Sobol, 2001)

O procedimento deverá distinguir intervalos que representam a
distribuição dos resultados entre execuções dos intervalos de confiança
relativos à estimativa de sua média. Essa distinção evita interpretar a
faixa de trajetórias como erro de medição ou como incerteza sobre os
casos reais de dengue. Para resultados acumulados, é importante definir
explicitamente o denominador populacional, a duração e a unidade
temporal, garantindo comparações consistentes entre cenários. (Gamal et
al., 2025; Morris, 1991; Sobol, 2001)

\subsubsection{5.5.3 Análise da sensibilidade espacial e temporal:
extensão metodológica
proposta}\label{anuxe1lise-da-sensibilidade-espacial-e-temporal-extensuxe3o-metodoluxf3gica-proposta}

Gamal et al. (2025) propõem a Spatio-Temporal Uncertainty Analysis
(ST-UA) para caracterizar a propagação, entre dois instantes, de medidas
de incerteza ou sensibilidade calculadas em unidades espaciais. O
trabalho aplica o método a um modelo de tabagismo, Tobacco Town, e
demonstra índices espaciais relacionados à relevância e à irrelevância
da propagação da incerteza, utilizando como exemplo um índice de efeitos
totais associado a Sobol. O presente simulador não reproduz esse modelo
nem seus parâmetros; a contribuição aproveitável é o desenho da análise,
cuja adaptação requer justificativa própria.

Uma aplicação inicial e factível ao município de São Paulo consistiria
em definir unidades territoriais não sobrepostas compatíveis com a base
GIS e registrar, para cada região i, tempo t e execução r, os
indicadores de infecções humanas incidentes, população de mosquitos por
estado e/ou presença de criadouros, conforme a disponibilidade e a
verificação dos registros. Para cada combinação (i,t), poderiam ser
produzidos mapas da média e do desvio-padrão entre replicações,
acompanhados por contagens ou taxas normalizadas pelo tamanho da
população sintética regional. Regiões com população sintética nula ou
reduzida deverão ser identificadas para evitar instabilidade e
interpretações enganosas. (Gamal et al., 2025; Morris, 1991; Sobol,
2001)

A comparação de mapas em dois ou mais instantes permitiria descrever a
evolução da dispersão espacial dos resultados, mas não constituiria
automaticamente a aplicação integral da ST-UA. A reprodução formal dos
índices do artigo exigiria definir previamente a medida de incerteza
U(i,t), justificar limiares e coeficientes dos índices, assegurar a
compatibilidade espacial e temporal e executar uma campanha de
amostragem apropriada. A mera diferença entre desvios-padrão de dois
dias não deverá ser denominada índice ST-UA. (Gamal et al., 2025;
Morris, 1991; Sobol, 2001)

Para avaliar os parâmetros de transmissão, abundância vetorial,
mobilidade e condições ambientais, propõe-se uma progressão: (a) triagem
de fatores com variações controladas; (b) análise de sensibilidade
global, se tecnicamente viável, com amostragem de combinações e
consideração das interações; (c) eventual estimação de índices de Sobol
e, em uma etapa avançada, sua associação à análise espaço-temporal. A
análise um-fator-por-vez é útil como triagem, mas não estima interações
entre fatores; índices de Sobol exigem delineamento amostral e custo
computacional substancialmente maiores. As milhares de execuções
utilizadas por Gamal et al. (2025) pertencem à aplicação daquele estudo
e não constituem uma exigência universal para este TCC.

Para manter a rastreabilidade, cada observação espacial deverá ser
vinculada ao identificador da execução, semente, versão do código,
parâmetros, data simulada, código da região e definição do indicador.
Antes da análise, serão verificadas a conservação ou a coerência das
contagens epidemiológicas, a estabilidade do relógio do modelo, o
sistema de coordenadas, a associação de agentes às regiões e a ausência
de duplicação dos eventos registrados. Quando essas condições não forem
satisfeitas, os mapas e índices deverão permanecer classificados como
não avaliados. (Oliveira et al., 2025; Garcias, 2024)

Quadro 5.5 --- Plano escalonado de análise da incerteza (não executado)

Etapa A --- Verificação: testar integridade dos estados, consistência
temporal e atribuição espacial; produto esperado: registros de teste e
correções documentadas. (HKIOSHI, 2026)

Etapa B --- Ensemble estocástico: replicar o cenário de referência com
sementes distintas; produto esperado: médias, dispersões, trajetórias e
precisão estimada. (Gamal et al., 2025; Morris, 1991; Sobol, 2001)

Etapa C --- Incerteza espaço-temporal descritiva: gerar indicadores por
região e por instante; produto esperado: mapas de médias, dispersões e
evolução temporal. (Gamal et al., 2025; Morris, 1991; Sobol, 2001)

Etapa D --- Sensibilidade global / ST-UA: somente se houver
instrumentação e orçamento computacional suficientes; produto esperado:
índices e mapas com delineamento e limitações explicitados. (Gamal et
al., 2025; Morris, 1991; Sobol, 2001)

Fonte: elaboração própria (2026), adaptado conceitualmente de Gamal et
al. (2025). Etapas propostas, sem resultados experimentais nesta versão.

\subsection{5.6 Configurações batch incorporadas ao código e verificação
de
execução}\label{configurauxe7uxf5es-batch-incorporadas-ao-cuxf3digo-e-verificauxe7uxe3o-de-execuuxe7uxe3o}

O Main.gaml do instantâneo identificado declara quatro experimentos
batch, apresentados no Quadro 5.6. As contagens de repetições e as
faixas de parâmetros abaixo descrevem a configuração encontrada no
código, não a quantidade de execuções bem-sucedidas comprovadas.
(HKIOSHI, 2026)

Quadro 5.6 --- Experimentos em lote identificados no código-fonte

\begin{longtable}[]{@{}
  >{\raggedright\arraybackslash}p{(\linewidth - 4\tabcolsep) * \real{0.3333}}
  >{\raggedright\arraybackslash}p{(\linewidth - 4\tabcolsep) * \real{0.3333}}
  >{\raggedright\arraybackslash}p{(\linewidth - 4\tabcolsep) * \real{0.3333}}@{}}
\toprule\noalign{}
\begin{minipage}[b]{\linewidth}\raggedright
\textbf{Identificador}
\end{minipage} & \begin{minipage}[b]{\linewidth}\raggedright
\textbf{Configuração declarada}
\end{minipage} & \begin{minipage}[b]{\linewidth}\raggedright
\textbf{Questão de verificação}
\end{minipage} \\
\midrule\noalign{}
\endhead
\bottomrule\noalign{}
\endlastfoot
E03Variabilidade & 30 repetições, cenário sem vacinação & Confirmar
sementes independentes e saídas por repetição. \\
E04Vacinacao & Cobertura de 0\% a 100\% em intervalos de 25\%; eficácia
de 100\%; 30 repetições & Confirmar que a cobertura modifica o
comportamento e que a eficácia é hipótese de cenário, não estimativa
clínica. \\
E05PopulacaoVetorial & 50 a 400 mosquitos, passos de 50; 30 repetições &
Confirmar vínculo do parâmetro com a criação efetiva de mosquitos. \\
E07Escalabilidade & 500 a 10.000 humanos, passos de 500; uma execução
por configuração & Confirmar tamanho populacional real, hardware, tempo
e memória medidos. \\
\end{longtable}

Fonte: elaboração própria, a partir das declarações do Main.gaml no
commit f3b3a424 (consulta de 30 out. 2026). Parâmetros e repetições
configurados; não representam resultados experimentais verificados.

\subsubsection{5.6.1 Dependências que condicionam a
interpretação}\label{dependuxeancias-que-condicionam-a-interpretauxe7uxe3o}

A inicialização dos mosquitos apresenta criação fixa de 100 indivíduos
em trecho examinado do código, enquanto o experimento E05 declara
variação de quantidade\_mosquito. Se o valor declarado não for conectado
ao comando de criação, diferentes níveis do experimento poderão gerar a
mesma condição inicial. Analogamente, a atribuição de uma semente fixa
no procedimento de inicialização pode comprometer a independência das
repetições. Esses comportamentos devem ser comprovados mediante
rastreamento do valor efetivo das variáveis por execução. (HKIOSHI,
2026)

Registro da auditoria histórica (versão de outubro de 2026): Também se
observou uma transição vetorial circular envolvendo a operação módulo 3
e uma inicialização com rnd(0,3), apesar de os indicadores
contabilizarem três estados. O teste deve enumerar as transições
permitidas e conferir a distribuição inicial, sem antecipar a causa
definitiva das divergências do CSV. O limite de ciclos deve ser
interpretado de acordo com o relógio efetivo, incluindo eventuais
multiplicadores de velocidade, antes de convertê-lo em dias
epidemiológicos. Observação: a issue \#12 relata alterações posteriores
em 7/11/2026; confirmar a situação atual no código e por testes antes de
repetir esta crítica como defeito vigente. (HKIOSHI, 2026)

\subsubsection{5.6.2 Critérios mínimos para considerar um experimento
concluído}\label{crituxe9rios-muxednimos-para-considerar-um-experimento-concluuxeddo}

Cada execução deverá registrar ID do experimento, ID da repetição,
semente efetiva, hash do commit, versão do GAMA, parâmetros resolvidos,
tamanho inicial e final das populações, duração simulada, número de
ciclos, tempo real de execução e localização dos arquivos de saída. Em
cada instante devem ser testadas as invariantes de contagem; falhas não
devem ser silenciosamente eliminadas da análise. Para E07, uma única
execução por tamanho não permite estimar adequadamente a dispersão do
desempenho, motivo pelo qual se recomendam repetições controladas.
(Gamal et al., 2025; Morris, 1991; Sobol, 2001)

\subsubsection{5.6.3 Interpretação preliminar dos registros, sem
atribuição de
cenário}\label{interpretauxe7uxe3o-preliminar-dos-registros-sem-atribuiuxe7uxe3o-de-cenuxe1rio}

No primeiro segmento do CSV do commit f3b3a424, o dia rotulado 0
registra 1.191 humanos e 100 mosquitos; no dia rotulado 30, o campo
humanos\_infectados atinge 382. A população vetorial cresce e decresce
ao longo do mesmo segmento. Essas observações documentam que o sistema
exportou séries dinâmicas, mas não permitem atribuir causalidade às
intervenções, inferir um efeito da vacinação ou estimar parâmetros
epidemiológicos reais sem identificar a configuração e a repetição de
origem. (HKIOSHI, 2026)

\subsection{5.7 Comparação com observações reais de São
Paulo}\label{comparauxe7uxe3o-com-observauxe7uxf5es-reais-de-suxe3o-paulo}

A avaliação epidemiológica poderá utilizar séries oficiais de casos do
município de São Paulo, preferencialmente desagregadas por semana
epidemiológica e, quando disponível, por distrito ou outra unidade
territorial compatível. O ano de 2024 constitui uma possível referência
por sua relevância epidemiológica, mas sua escolha deverá ser
acompanhada de justificativa do período, população, parâmetros,
condições ambientais e dados de inicialização utilizados no simulador.
(Oliveira et al., 2025; Garcias, 2024)

É necessário estabelecer a distinção entre infecções internas do modelo
e casos observados. A notificação real depende de manifestações
clínicas, acesso ao diagnóstico e processos de vigilância, ao passo que
o contador de infecções simuladas pode representar eventos diferentes. A
comparação poderá exigir um submodelo de observação, reamostragem
temporal e indicadores normalizados, sem confundir uma população
sintética reduzida com a população municipal real. (HKIOSHI, 2026)

Somente após a construção de séries comparáveis poderá ser estimada uma
medida de erro, como RMSE = raiz{[} (1/n) × soma\_t (y\_sim,t −
y\_obs,t)\^{}2 {]}, ou outras métricas predefinidas. Ajustar parâmetros
e avaliar o modelo com a mesma série, sem separação de dados, enfraquece
a evidência de generalização. A validade deverá ser apresentada de
maneira proporcional às bases e aos testes efetivamente realizados.
(HKIOSHI, 2026)

\subsection{5.8 Avaliação do desempenho
computacional}\label{avaliauxe7uxe3o-do-desempenho-computacional}

A avaliação de desempenho, componente central para um TCC de Ciência da
Computação, deverá caracterizar o custo de inicialização, o
processamento por ciclo e a memória consumida, controlando o número de
agentes, as configurações de visualização e a resolução espacial. Os
testes deverão ser feitos em hardware identificado, com aquecimento e
execução padronizados quando aplicável, distinguindo diferenças de
desempenho das diferenças de comportamento estocástico. (HKIOSHI, 2026)

Quadro 5.4 --- Cenários experimentais de desempenho (planejados)

\begin{longtable}[]{@{}
  >{\raggedright\arraybackslash}p{(\linewidth - 4\tabcolsep) * \real{0.3333}}
  >{\raggedright\arraybackslash}p{(\linewidth - 4\tabcolsep) * \real{0.3333}}
  >{\raggedright\arraybackslash}p{(\linewidth - 4\tabcolsep) * \real{0.3333}}@{}}
\toprule\noalign{}
\begin{minipage}[b]{\linewidth}\raggedright
\textbf{Fator}
\end{minipage} & \begin{minipage}[b]{\linewidth}\raggedright
\textbf{Variação proposta}
\end{minipage} & \begin{minipage}[b]{\linewidth}\raggedright
\textbf{Métrica}
\end{minipage} \\
\midrule\noalign{}
\endhead
\bottomrule\noalign{}
\endlastfoot
População humana & Configurações crescentes e documentadas & Tempo total
e por ciclo \\
População vetorial & Configurações crescentes e documentadas & Custo das
interações e memória \\
Visualização & Com e sem renderização gráfica, quando suportado &
Sobrecarga de interface \\
Resolução territorial & Camadas e agregações comparáveis & Tempo de
carga e processamento \\
Consultas espaciais & Estratégia implementada e alternativa, se viável &
Tempo por consulta/ciclo \\
\end{longtable}

Fonte: elaboração própria (2026). Experimentos não executados nesta
revisão.

Sem medições reais, não é possível afirmar escalabilidade, superioridade
computacional em relação a NetLogo, Mesa ou Repast, nem complexidade
prática garantida das consultas espaciais. Caso implementadas por
varredura completa dos agentes humanos e mosquitos, comparações par a
par podem crescer proporcionalmente ao produto das populações;
entretanto, a avaliação da implementação GAMA requer medidas e inspeção
das operações realmente chamadas. (Oliveira et al., 2025; Garcias, 2024)

\subsection{5.9 Discussão crítica e
limitações}\label{discussuxe3o-cruxedtica-e-limitauxe7uxf5es}

O arquivo consultado demonstra que o software dispõe de uma rotina de
registro de séries temporais, com variação dos estados humanos e
vetoriais no intervalo exportado. Entretanto, as diferenças entre totais
e estados, o crescimento de uma unidade na população humana e a ausência
de metadados suficientes impedem classificar essas saídas como validação
do modelo. Tais limitações são especialmente relevantes porque um
simulador pode produzir curvas visualmente plausíveis apesar de conter
erros de implementação ou hipóteses inadequadas. (HKIOSHI, 2026)

Os trabalhos de Oliveira et al. (2025) e Garcias (2024) oferecem
contexto para a modelagem da dengue em GAMA, enquanto Gamal et al.
(2025) fundamentam a distinção entre variabilidade estocástica,
incerteza paramétrica e análise espaço-temporal. Seus resultados
numéricos e cenários são externos a este projeto; não constituem
validação automática das saídas do TccAbm.

Portanto, a contribuição computacional demonstrável nesta etapa está na
análise rastreável do software, na descrição das evidências existentes e
na especificação de um protocolo capaz de produzir resultados
auditáveis. A confirmação experimental da correção do código, da
adequação epidemiológica e do desempenho ainda requer novas execuções e
documentação sistemática. (HKIOSHI, 2026)

A aplicação de métodos de incerteza espaço-temporal, inspirados em Gamal
et al. (2025), somente será considerada após a instrumentação de
resultados desagregados por território e o controle de replicações. A
presença de heterogeneidade espacial ou de dispersão entre sementes, por
si só, não valida a representação epidemiológica nem estabelece
capacidade preditiva.

\subsection{5.10 Síntese do capítulo}\label{suxedntese-do-capuxedtulo-3}

Foi identificado um arquivo CSV com 229 registros e quatro experimentos
em lote declarados no código-fonte do instantâneo analisado. A inspeção
exploratória evidenciou divergências na contabilização dos estados
vetoriais, ausência de identificadores adequados de execução e pontos
que exigem teste no controle temporal, nas sementes e na parametrização.
O capítulo distingue evidências registradas, configurações implementadas
e etapas ainda não executadas ou não verificadas. Antes de concluir
sobre sensibilidade, vacinação, escalabilidade ou validade
epidemiológica, são necessárias execuções controladas e reproduzíveis.
(Gamal et al., 2025; Morris, 1991; Sobol, 2001)

\section{6 CONCLUSÃO E TRABALHOS
FUTUROS}\label{conclusuxe3o-e-trabalhos-futuros}

\subsection{6.1 Síntese da pesquisa}\label{suxedntese-da-pesquisa}

O presente trabalho abordou a especificação, a documentação e a
avaliação computacional em desenvolvimento de um simulador espacial e
estocástico de transmissão da dengue, implementado na plataforma GAMA e
delimitado ao município de São Paulo. O problema central investigado foi
a construção de uma representação baseada em agentes capaz de articular
entidades humanas, mosquitos e ambiente geográfico, com comportamentos
individuais e mecanismos probabilísticos de interação. Ao longo do
texto, foram caracterizados os conceitos epidemiológicos necessários, a
integração entre dados GIS e simulação, a estrutura modular do código
GAML e os critérios indispensáveis para uma avaliação experimental
reproduzível. (Gamal et al., 2025)

A contribuição computacional do estudo consiste no exame estruturado da
arquitetura e do comportamento implementado, na formalização dos
requisitos experimentais e na explicitação dos limites da evidência
disponível. A organização dos agentes e das regras de interação oferece
um objeto concreto para discutir simulação baseada em agentes, controle
temporal, estados epidemiológicos, parâmetros, persistência de saídas e
integridade de indicadores. O trabalho não assume, por essa razão, que
representar computacionalmente um mecanismo epidemiológico seja
equivalente a demonstrar sua fidelidade à dinâmica observada da dengue.

\subsection{6.2 Contribuições documentadas e estado das
verificações}\label{contribuiuxe7uxf5es-documentadas-e-estado-das-verificauxe7uxf5es}

A análise dos arquivos do repositório identificou a existência de
módulos de inicialização, espécies de agentes humanos e vetoriais,
elementos ambientais e configurações de experimentos em lote. O código
disponível foi organizado para permitir cenários de variabilidade
estocástica, vacinação, população vetorial e escalabilidade. No
acompanhamento de 7 de novembro de 2026, commits vinculados à issue \#12
registraram alterações destinadas a corrigir inconsistências
relacionadas a sementes, quantidade inicial de mosquitos, inicialização
dos estados vetoriais e armazenamento dos resultados. Tais registros
evidenciam evolução do projeto e revisão orientada por testes, mas não
comprovam, isoladamente, a execução satisfatória dos cenários. (Gamal et
al., 2025)

O conjunto histórico de dados examinado na revisão de outubro permitiu
identificar uma divergência entre o total de mosquitos e a soma das
categorias epidemiológicas registradas no primeiro instante observado.
Essa divergência teve valor metodológico como indicação da necessidade
de estabelecer invariantes de população e testes automatizados. A
confirmação da correção dessa condição e dos demais critérios de
aceitação da issue \#12 depende de executar o código atualizado e
confrontar as saídas por cenário, réplica e semente. (Taillandier et
al., 2019; Macal; North, 2010)

\subsection{6.3 Limitações}\label{limitauxe7uxf5es}

A principal limitação desta versão consolidada é a ausência, nos
materiais analisados para sua preparação, de um pacote comprovadamente
completo de saídas experimentais associadas aos commits atuais, contendo
configuração, semente efetiva, duração simulada e série temporal para
cada execução. Consequentemente, não é metodologicamente adequado
afirmar diferenças estatisticamente demonstradas entre cenários de
vacinação, variação vetorial e escalabilidade. Também permanecem
pendentes a confirmação integral das fontes e dos tratamentos de dados
climáticos e geográficos, a especificação operacional de parâmetros
epidemiológicos e a verificação de consistência temporal entre ciclos e
dias.

O emprego de uma população sintética reduzida, a simplificação dos
estados e das condições de transmissão e a inexistência de um
procedimento de calibração e avaliação externa completo limitam a
interpretação epidemiológica dos resultados. As saídas do simulador não
podem ser convertidas diretamente em previsão do número de casos
notificados no município de São Paulo, porque diferem da observação real
em tamanho populacional, cobertura espacial efetiva, subnotificação e
regras de observação. Essas limitações são particularmente relevantes em
modelos explicativos e exploratórios, cujas conclusões devem permanecer
condicionadas às hipóteses representadas.

\subsection{6.4 Trabalhos futuros e critérios para conclusão
experimental}\label{trabalhos-futuros-e-crituxe9rios-para-conclusuxe3o-experimental}

A primeira prioridade é validar os critérios de aceite associados à
correção de código: executar o cenário E03 com sementes distintas e
confirmar séries distintas; repetir uma semente e confirmar
reprodutibilidade; verificar a igualdade entre o total de mosquitos e a
soma de seus estados; e assegurar a rastreabilidade por cenário, réplica
e semente. Em seguida, é necessário executar os cenários E04 e E05,
confirmar que os parâmetros manipulados afetam efetivamente a
inicialização ou os eventos do modelo, e construir tabelas de
estatísticas descritivas sobre replicações independentes. A avaliação de
escalabilidade deverá utilizar configurações controladas de ambiente e
registrar tempo de execução e consumo de recursos, sem inferir
desempenho a partir de um único teste. (Taillandier et al., 2019; Macal;
North, 2010)

Como extensões científicas, recomenda-se investigar a sensibilidade
global dos parâmetros após a estabilização dos resultados, aprimorar a
representação da heterogeneidade espacial e comparar indicadores
adequadamente normalizados com séries oficiais de dengue no município. A
aplicação de técnicas de quantificação de incerteza espaço-temporal,
como discutido por Gamal et al. (2025), requer saídas por unidade
geográfica, múltiplas execuções e controle rigoroso das fontes de
variabilidade. Somente depois de cumpridas essas condições será possível
avaliar a contribuição de técnicas de calibração, análise de
sensibilidade global e eventual previsão de cenários.

\subsection{6.5 Considerações finais}\label{considerauxe7uxf5es-finais}

O trabalho demonstra a pertinência da modelagem baseada em agentes e da
integração GIS para a investigação computacional de mecanismos
associados à transmissão da dengue. Demonstra igualmente que um
simulador cientificamente defensável depende não apenas de uma
arquitetura funcional, mas de verificações de consistência, controle dos
experimentos, documentação dos parâmetros e interpretação responsável
das saídas. A conclusão definitiva sobre os objetivos experimentais
permanecerá condicionada à apresentação das evidências exigidas pelo
orientador, sem substituir resultados ainda não obtidos por descrições
de intenção. (Taillandier et al., 2019; Macal; North, 2010)

\section{REFERÊNCIAS}\label{referuxeancias}

ANYLOGIC. AnyLogic Help. {[}S. l.{]}: AnyLogic, {[}s. d.{]}. Disponível
em: https://anylogic.help/. Acesso em: 8 nov. 2026.

ARIFIN, S. M. Niaz; MADEY, Gregory R.; COLLINS, Frank H. Spatial
agent-based simulation modeling in public health: design,
implementation, and applications for malaria epidemiology. Hoboken: John
Wiley \& Sons, 2016. DOI: https://doi.org/10.1002/9781118964385.

BANOS, Arnaud; LANG, Christophe; MARILLEAU, Nicolas (org.). Agent-based
spatial simulation with NetLogo: volume 1: introduction and bases.
London: ISTE Press; Elsevier, 2015. ISBN 9781785480553.

BRASIL. Ministério da Saúde. Dengue. Brasília, DF: Ministério da Saúde,
{[}s. d.{]}. Disponível em:
https://www.gov.br/saude/pt-br/assuntos/saude-de-a-a-z/d/dengue. Acesso
em: 8 nov. 2026.

DOEGE, Daniel Garbin. Avaliação de estratégias de controle da dengue por
modelagem baseada em agentes. 2026. Trabalho de Conclusão de Curso
(Graduação em Biomedicina) --- Universidade Federal de Ciências da Saúde
de Porto Alegre, Porto Alegre, 2026. Disponível em:
https://repositorio.ufcspa.edu.br/handle/123456789/3553. Acesso em: 8
nov. 2026.

GAMA PLATFORM. GAMA Platform: documentação oficial. {[}S. l.{]}: GAMA
Platform, {[}s. d.{]}. Disponível em: https://gama-platform.org/. Acesso
em: 8 nov. 2026.

GAMAL, Yahya et al. An analysis of spatial and temporal uncertainty
propagation in agent-based models. Philosophical Transactions of the
Royal Society A: Mathematical, Physical and Engineering Sciences,
London, v. 383, n. 2293, art. 20240229, 2025. DOI:
https://doi.org/10.1098/rsta.2024.0229.

GARCIAS, David Batista. Modelo de simulação da transmissão da dengue
baseado em agentes. 2024. Dissertação (Mestrado) --- Universidade NOVA
de Lisboa, Portugal, 2024. {[}Conferir unidade acadêmica e dados
catalográficos na folha de rosto da dissertação{]}.

GHARAKHANLOU, N. M.; HOOSHANGI, N.; HELBICH, M. A spatial agent-based
model to assess the spread of malaria in relation to anti-malaria
interventions in Southeast Iran. ISPRS International Journal of
Geo-Information, {[}S. l.{]}, 2020. {[}Conferir volume, fascículo e
identificador do artigo{]}.

GHARAKHANLOU, N. M.; MESGARI, M. S.; HOOSHANGI, N. Developing an
agent-based model for simulating the dynamic spread of Plasmodium vivax
malaria: a case study of Sarbaz, Iran. Ecological Informatics, {[}S.
l.{]}, v. 54, 2019. {[}Conferir identificador do artigo e DOI{]}.

GRIMM, Volker et al. The ODD protocol for describing agent-based and
other simulation models: a second update to improve clarity,
replication, and structural realism. Journal of Artificial Societies and
Social Simulation, {[}S. l.{]}, v. 23, n. 2, art. 7, 2020. DOI:
https://doi.org/10.18564/jasss.4259.

HINSCH, Martin. Introduction to modelling: part 1. YouTube, {[}s. d.{]}.
Disponível em: https://www.youtube.com/watch?v=V5XyzCQa6uo. Acesso em:
27 out. 2026. {[}Confirmar título e data diretamente no vídeo{]}.

HKIOSHI. TccAbm: repositório de código-fonte e dados do simulador.
GitHub, 2026. Disponível em: https://github.com/hkioshi/TccAbm. Acesso
em: 8 nov. 2026. Para observações históricas, consultar o commit
específico indicado no corpo do texto.

ISIDORO, Carlos; FACHADA, Nuno; BARATA, Fábio; ROSA, Agostinho.
Agent-based model of dengue disease transmission by Aedes aegypti
populations. In: ADVANCES IN ARTIFICIAL LIFE: DARWIN MEETS VON NEUMANN.
Berlin: Springer, 2011. p. 345--352. DOI:
https://doi.org/10.1007/978-3-642-21283-3\_43.

JINDAL, Akshay; RAO, Shrisha. Agent-based modeling and simulation of
mosquito-borne disease transmission. In: INTERNATIONAL CONFERENCE ON
AUTONOMOUS AGENTS AND MULTIAGENT SYSTEMS, 16., 2017, São Paulo.
Proceedings {[}...{]}. {[}S. l.{]}: International Foundation for
Autonomous Agents and Multiagent Systems, 2017. p. 426--435. Disponível
em: https://aamas.csc.liv.ac.uk/Proceedings/aamas2017/pdfs/p426.pdf.
Acesso em: 8 nov. 2026.

MACAL, Charles M.; NORTH, Michael J. Tutorial on agent-based modelling
and simulation. Journal of Simulation, {[}S. l.{]}, v. 4, n. 3, p.
151--162, 2010. DOI: https://doi.org/10.1057/jos.2010.3.

MESA-GEO. Mesa-Geo: GIS extension for Mesa. {[}S. l.{]}: Mesa-Geo, {[}s.
d.{]}. Disponível em: https://mesa-geo.readthedocs.io/. Acesso em: 8
nov. 2026.

MESA. Mesa: agent-based modeling in Python. {[}S. l.{]}: Mesa, {[}s.
d.{]}. Disponível em: https://mesa.readthedocs.io/. Acesso em: 8 nov.
2026.

MORRIS, Max D. Factorial sampling plans for preliminary computational
experiments. Technometrics, {[}S. l.{]}, v. 33, n. 2, p. 161--174, 1991.
DOI: https://doi.org/10.1080/00401706.1991.10484804.

NETLOGO. GIS Extension: NetLogo User Manual. {[}S. l.{]}: NetLogo, {[}s.
d.{]}. Disponível em: https://docs.netlogo.org/gis.html. Acesso em: 8
nov. 2026.

OLIVEIRA, Emily Brito de et al. Modelo de transmissão da dengue por
simulação multi-agentes. In: SIMPÓSIO BRASILEIRO DE PESQUISA
OPERACIONAL, 57., 2025, Gramado. Anais {[}...{]}. {[}S. l.{]}: Galoá,
2025. DOI: https://doi.org/10.59254/sbpo-2025-212362.

REPAST. Repast Simphony: documentação. {[}S. l.{]}: Repast, {[}s. d.{]}.
Disponível em: https://repast.github.io/. Acesso em: 8 nov. 2026.

SANDERSON, Grant. Simulating an epidemic. 3Blue1Brown, 27 mar. 2020.
Vídeo (23 min). Disponível em:
https://www.youtube.com/watch?v=gxAaO2rsdIs. Acesso em: 27 out. 2026.

SANTOS, Aline Rodrigues; SANTOS, Fernando. Introductory guide to
agent-based simulation development on the GAMA platform. In:
WORKSHOP-ESCOLA DE SISTEMAS DE AGENTES, SEUS AMBIENTES E APLICAÇÕES
(WESAAC), 18., 2024, Brasília. Anais {[}...{]}. Porto Alegre: Sociedade
Brasileira de Computação, 2024. p. 96--107. DOI:
https://doi.org/10.5753/wesaac.2024.33459.

SOBOL, Ilya M. Global sensitivity indices for nonlinear mathematical
models and their Monte Carlo estimates. Mathematics and Computers in
Simulation, {[}S. l.{]}, v. 55, n. 1--3, p. 271--280, 2001. DOI:
https://doi.org/10.1016/S0378-4754(00)00270-6.

TAILLANDIER, Patrick et al. Building, composing and experimenting
complex spatial models with the GAMA platform. GeoInformatica, {[}S.
l.{]}, v. 23, p. 299--322, 2019. DOI:
https://doi.org/10.1007/s10707-018-00339-6.

TAILLANDIER, Patrick; DROGOUL, Alexis. From GIS data to GIS agents,
modeling with the GAMA simulation platform. In: TECHNICAL FORUM GROUP ON
AGENT AND MULTI-AGENT-BASED SIMULATION, 2010. {[}Anais...{]}. {[}S. l.:
s. n.{]}, 2010. {[}Conferir os dados do evento no documento original{]}.

WORLD HEALTH ORGANIZATION. Dengue. Geneva: WHO, 21 ago. 2025. Disponível
em:
https://www.who.int/news-room/fact-sheets/detail/dengue-and-severe-dengue.
Acesso em: 8 nov. 2026.

Nota técnica bibliográfica: referências normalizadas no padrão
autor-data da ABNT. As entradas de Garcias (2024), Gharakhanlou et al.
(2019/2020), Hinsch ({[}s. d.{]}) e Taillandier e Drogoul (2010) ainda
requerem conferência catalográfica final nos documentos originais,
indicada entre colchetes; essas indicações não devem permanecer na
versão de depósito. Esta revisão não constitui validação científica do
simulador e não substitui a compilação da bibliografia no LaTeX
institucional.

\section{APÊNDICE A --- RASTREABILIDADE DAS EXIGÊNCIAS DO
ORIENTADOR}\label{apuxeandice-a-rastreabilidade-das-exiguxeancias-do-orientador}

\begin{longtable}[]{@{}
  >{\raggedright\arraybackslash}p{(\linewidth - 4\tabcolsep) * \real{0.3333}}
  >{\raggedright\arraybackslash}p{(\linewidth - 4\tabcolsep) * \real{0.3333}}
  >{\raggedright\arraybackslash}p{(\linewidth - 4\tabcolsep) * \real{0.3333}}@{}}
\toprule\noalign{}
\begin{minipage}[b]{\linewidth}\raggedright
\textbf{Exigência / issue}
\end{minipage} & \begin{minipage}[b]{\linewidth}\raggedright
\textbf{Providência textual nesta versão}
\end{minipage} & \begin{minipage}[b]{\linewidth}\raggedright
\textbf{Evidência ainda necessária}
\end{minipage} \\
\midrule\noalign{}
\endhead
\bottomrule\noalign{}
\endlastfoot
\#13 Resumo/Abstract & Reescritos para GAMA/GAML, sem Mesa, ML ou
validação não demonstrada & Atualizar números somente após resultados e
recompilar LaTeX/BibTeX \\
\#8 Metodologia & Capítulos 3, 4 e 5 especificam agentes, parâmetros,
verificação e cenários & Confirmar fonte, ano, CRS, clima, valores e
ciclos no código \\
\#12 Correção experimental & Nota de atualização registra commits 7/11 e
critérios de aceitação & E03 por semente, E04 eficácia 50\%, E05
população, estados S+E+I \\
\#11 Resultados/Conclusão & Capítulo 5 distingue dados históricos de
resultados futuros; capítulo 6 limita inferências & CSV rastreáveis e
gráficos com replicações \\
\#16 Prazo 12/11 & Texto consolidado e roteiro documental & Publicar no
main, PDF compilado e comentários nas issues \\
\end{longtable}

\subsection{A.1 Protocolo mínimo de validação das
execuções}\label{a.1-protocolo-muxednimo-de-validauxe7uxe3o-das-execuuxe7uxf5es}

1. Registrar commit SHA do código, versão do GAMA, data, sistema
operacional, configuração da máquina e fontes de dados.

2. Executar E03 com duas sementes distintas; comparar séries e comprovar
diferenças; repetir uma das sementes para comprovar igualdade sob o
mesmo ambiente.

3. Confirmar em cada CSV as colunas cenário, réplica, semente, ciclo,
dia acumulado, totais e compartimentos epidemiológicos.

4. Garantir invariantes: total de mosquitos = S + E + I, sob definição
de categorias mutuamente exclusivas; mesmo princípio para humanos
conforme estados de recuperação, morte e outros estados presentes.

5. No E05, verificar a correspondência entre quantidade\_mosquito
configurada e a contagem inicial registrada.

6. No E04, testar eficácia 50\% em cenário controlado e comparar a
probabilidade de transmissão efetiva na escala 0--100, sem exigir que
toda réplica apresente infecção de vacinados.

7. Calcular média, desvio-padrão, mediana, percentis e total de novas
infecções por cenário, somente após verificar consistência e número de
réplicas.

8. Publicar CSVs, protocolo, gráficos e PDF; comentar cada issue com
commit, caminho dos arquivos e evidências reproduzíveis.

\subsection{A.2 Limitações desta
entrega}\label{a.2-limitauxe7uxf5es-desta-entrega}

Nota de incorporação bibliográfica (8 nov. 2026): as referências de
Arifin, Madey e Collins (2016), Banos, Lang e Marilleau (2015), Isidoro
et al. (2011), Jindal e Rao (2017) e Doege (2026) foram acrescentadas às
seções metodológicas ou comparativas pertinentes. Não foram utilizados
resultados desses autores como se fossem produzidos pelo simulador GAMA
do presente projeto. Permanecem necessárias a execução dos experimentos,
a conferência bibliográfica integral e a integração ao template LaTeX
institucional.

Este DOCX é uma consolidação editorial de versões prévias dos capítulos
disponíveis nesta conversa, acrescida de correções conceituais e de uma
nota histórica sobre o estado mais recente observado nas issues. Não
substitui a verificação da árvore completa dos arquivos do repositório,
não altera o GitHub, não compila o projeto LaTeX e não executa o modelo
GAMA. Por isso, o atendimento integral às exigências do orientador
permanece condicionado às ações técnicas descritas neste apêndice.
